% ============================================================
\section{Traffic System Analysis: A Meso-Micro Simulation Framework
         for Mixed Fuel and Electric Vehicles}
\label{sec:traffic_system_analysis_ltm}
% ============================================================

\begin{remark}[Implementation target]
  This section is the mathematical contract for the
  \texttt{hacdcpf::evpt} traffic sub-module.  Function signatures,
  data structures, and algorithmic steps given here are the specification
  that the C++ implementation in \texttt{src/ev\_power\_traffic/} follows.
  Where a function is already implemented a cross-reference to the source
  file is given; where it is planned, the label \planned{} appears.
\end{remark}

% ------------------------------------------------------------
\subsection{Model Architecture}
\label{sec:traffic-architecture}
% ------------------------------------------------------------

The framework is organised as a five-layer pipeline:

\begin{equation}
  \underbrace{\text{Vehicle behaviour}}_{\text{micro}}
  \;\xrightarrow{\text{aggregate}}\;
  \underbrace{\text{Fleet flows}}_{\text{meso input}}
  \;\xrightarrow{\text{DNL}}\;
  \underbrace{\text{Traffic propagation}}_{\text{CTM / LTM}}
  \;\xrightarrow{\text{service}}\;
  \underbrace{\text{Station dynamics}}_{\text{Q, B, H}}
  \;\xrightarrow{\text{feedback}}\;
  \underbrace{\text{Equilibrium / optimisation}}_{\text{UE or SO}}
  \label{eq:pipeline}
\end{equation}

\noindent
Each layer has a precisely defined input--output contract:
\begin{itemize}
  \item \emph{Vehicle behaviour layer}: answers how individual vehicles
        choose routes and charging plans given perceived travel times,
        queue lengths, and energy prices.
  \item \emph{Fleet aggregation layer}: converts large populations of
        identically-behaving vehicles into continuous OD flows, making
        Dynamic Network Loading (DNL) tractable.
  \item \emph{CTM/LTM layer}: propagates fleet flows through the road
        network, computing shockwaves, spillback, travel times, and
        FIFO-consistent trajectories.
  \item \emph{Station service layer}: models charging-station queues,
        plug utilisation, service completions, and dwell populations as
        endogenous nodes inside the network loading.
  \item \emph{Equilibrium/optimisation layer}: closes the loop;
        at equilibrium, no vehicle has an incentive to deviate
        (Wardrop~I for ICVs, economic UE for EVs).
\end{itemize}

\noindent\textbf{This section covers:}
\begin{itemize}
  \item micro ICV and EV behaviour models (route choice, fuel consumption,
        SOC dynamics, charging and V2G decisions);
  \item fleet aggregation from micro behaviour to OD flow variables;
  \item CTM-based system-level network loading with cell-wise spatial
        state, route-commodity tracking, and DUE via MSA or variational
        inequality methods; system-optimal dynamic traffic assignment is
        formulated as an LP under triangular-FD assumptions;
  \item LTM-based meso-micro loading with cumulative curves, FIFO
        trajectory reconstruction, and an individual EV interaction layer;
  \item charging stations as endogenous service nodes;
  \item the dynamic user-equilibrium fixed-point problem and MSA solution;
  \item system-optimal LP, MILP infrastructure design, and MPC extensions.
\end{itemize}

\noindent\textbf{This section does not cover:}
\begin{itemize}
  \item optimal power flow or grid dispatch (Section~\ref{sec:ev_power_traffic_dynamic_model});
  \item carbon-flow or resilience analysis;
  \item full agent-level simulation of every vehicle.
\end{itemize}

% ------------------------------------------------------------
\subsection{Basic Sets, Indices, and Network Definition}
\label{sec:traffic-network}
% ------------------------------------------------------------

\subsubsection{Time horizon}

The discrete time set is $\mathcal{K} = \{0, 1, \ldots, T\}$ with
uniform step size $\Delta t$ (hours).  A flow expressed in veh/h is
multiplied by $\Delta t$ to obtain vehicles per step.

\subsubsection{Road network}

The road network is a directed graph
$\mathcal{G}^{R} = (\mathcal{N}^{R}, \mathcal{A}^{R})$,
where $\mathcal{N}^{R}$ is the set of road nodes and
$\mathcal{A}^{R}$ the set of directed road links.
Each link $a \in \mathcal{A}^{R}$ is characterised by:
\begin{align*}
  L_a         &\quad \text{length (km)}, \\
  v^{ff}_a    &\quad \text{free-flow speed (km/h)}, \\
  w_a         &\quad \text{backward wave speed (km/h)}, \\
  \bar{q}_a   &\quad \text{capacity (veh/h)}, \\
  N^{jam}_a   &\quad \text{jam storage (veh)}.
\end{align*}

The triangular fundamental diagram (FD) is
\begin{equation}
  q_a(k_a) =
  \begin{cases}
    v^{ff}_a \, k_a           & k_a \le k^c_a, \\
    w_a(k^j_a - k_a)          & k_a > k^c_a,
  \end{cases}
  \label{eq:triangular-fd}
\end{equation}
with critical density $k^c_a = \bar{q}_a / v^{ff}_a$ and jam density
$k^j_a = \bar{q}_a (1/v^{ff}_a + 1/w_a)$.
The backward wave speed satisfies
\begin{equation}
  w_a = \frac{v^{ff}_a \, k^c_a}{k^j_a - k^c_a}.
  \label{eq:wave-speed}
\end{equation}

\subsubsection{OD pairs, vehicle classes, and demand}
\label{sec:traffic-demand}

The OD set is $\mathcal{W}$; a pair $w = (o, d) \in \mathcal{W}$.
The vehicle class set is $\mathcal{C} = \{\mathrm{F}, \mathrm{E}\}$
(fuel vehicles and EVs).

\paragraph{Full DUE mode (departure time endogenous).}
Each OD pair $w$ has a fixed total demand $D_w$ over a departure window
$\mathcal{K}^{dep}_w \subseteq \mathcal{K}$ and a desired arrival time
$t^*_w$.  Departure time is an endogenous choice variable.  Given EV
penetration $p^{EV}_w \in [0,1]$:
\begin{equation}
  D^F_w = (1 - p^{EV}_w)\,D_w, \qquad
  D^E_w = p^{EV}_w\,D_w.
  \label{eq:demand-total}
\end{equation}
The departure profile $d_{w,k} = \sum_r h^F_{w,r,k} + \sum_\omega h^E_{w,\omega,k}$
is an equilibrium \emph{output}, not a given.  This is the full DUE
specification; the current CTM-DUE implementation supports this mode for
EV demands with non-empty \texttt{departure\_window\_steps}.

\paragraph{Fixed-departure mode (departure time exogenous).}
When an exogenous departure profile $d_{w,k}$ is supplied, the model
reduces to a fixed-departure dynamic assignment: equilibrium is imposed
only across routes or activity plans \emph{within} each departure step
$k$.  In this case the demand split is
\begin{equation}
  d^F_{w,k} = (1 - p^{EV}_{w,k})\,d_{w,k}, \qquad
  d^E_{w,k} = p^{EV}_{w,k}\,d_{w,k}.
  \label{eq:demand-split}
\end{equation}
This is a restricted special case of the full DUE; benchmark tests in
\S\ref{sec:traffic-numerical-results} use fixed-departure mode for
comparability with published results.

\subsubsection{Route and activity-plan sets}

The ICV route set for OD $w$ is $\mathcal{R}_w$.
An EV \emph{activity plan} $\omega = (r, s, \chi)$ combines a road
path $r$, a \emph{single} optional charging stop $s \in \mathcal{S} \cup \{\varnothing\}$,
and a charging/V2G strategy $\chi$ (charging power profile, charging energy,
discharging energy, dwell duration).
The current implementation supports single-stop plans.
Multi-stop plans (secondary charging after V2G discharge) are a planned
rolling-horizon extension; until then, secondary charging after V2G must
be handled by including a separate plan with the net energy already
accounting for the discharge.
The full activity-plan set for OD $w$ at step $k$ is $\Omega_{w,k}$.
ICVs choose a route; EVs choose an activity plan.

\begin{remark}[Activity-plan generation: theory vs.\ implementation]
  The description below gives the full mathematical specification.
  In the current implementation, $\Omega_{w,k}$ is \emph{not} generated
  automatically.  The calling code provides pre-specified route
  alternatives and charging-stop templates directly via
  \texttt{EVDemand::candidate\_route\_indices} and
  \texttt{RouteAlternative::charging\_stops}
  (\texttt{types.hpp}).  Automatic generation of $\Omega_{w,k}$ by
  $K$-shortest-path enumeration, station reachability, and
  charging-template synthesis is a planned extension and not yet
  implemented in \texttt{evpt\_internal.hpp}.
\end{remark}

Formally, $\Omega_{w,k}$ is defined as
\begin{equation}
  \Omega_{w,k}
  =
  \bigl\{
    (r,s,\chi) :
    r\in\mathcal{R}^K_w,\;
    s\in\mathcal{S}_r\cup\{\varnothing\},\;
    \chi\in\Xi_{s,k},\;
    e_{\ell}-E^{rem}_{\ell}\ge e^{res},\;\forall\ell
  \bigr\}.
  \label{eq:plan-set}
\end{equation}

\subsubsection{Charging stations}

The station set $\mathcal{S}$ collects service nodes embedded in the
road network.  Station $s$ is characterised by plug count $\bar{B}_s$,
maximum charging power $\bar{p}^{ch}_s$, maximum discharging power
$\bar{p}^{dis}_s$, queue capacity $\bar{Q}_s$, dwell capacity $\bar{H}_s$,
and grid-bus mapping $\mathrm{bus}(s)$.

% ------------------------------------------------------------
\subsection{Modelling Assumptions}
\label{sec:traffic-assumptions}
% ------------------------------------------------------------

\begin{description}
  \item[\textbf{A1} (LWR kinematic wave).] Road links follow first-order
    traffic flow theory (LWR) with triangular FD~\eqref{eq:triangular-fd}.
  \item[\textbf{A2} (FIFO).] Vehicles on the same link satisfy
    first-in-first-out: if $k^{in}_{i,a} \le k^{in}_{j,a}$ then
    $k^{out}_{i,a} \le k^{out}_{j,a}$.
  \item[\textbf{A3} (Equal PCE).] ICVs and EVs have identical
    passenger-car equivalence: $\mathrm{PCE}^F = \mathrm{PCE}^E = 1$.
    Class-specific PCE may be added as an extension.
  \item[\textbf{A4} (CTM/LTM for propagation).] CTM and LTM handle
    traffic propagation only; vehicle route choice and station activity
    are determined by the vehicle behaviour layer.
  \item[\textbf{A5} (ICV minimum-time).] The base ICV objective is
    minimum travel time.  Fuel costs may be included in the generalised
    cost as an extension.
  \item[\textbf{A6} (EV economic optimum).] EVs minimise a generalised
    cost that includes travel time, station waiting, charging service
    time, energy price, battery degradation, and SOC feasibility
    (Section~\ref{sec:traffic-ev-cost}).
  \item[\textbf{A7} (Fleet representation).] Vehicles sharing the same
    OD, departure step, and behaviour type are aggregated into a fleet
    flow $x^F_{w,r,k}$ or $x^E_{w,\omega,k}$ (Section~\ref{sec:traffic-fleet}).
  \item[\textbf{A8} (Discrete SOC dynamics).] EV battery energy follows
    the discrete update~\eqref{eq:soc-evolution}.
  \item[\textbf{A9} (Mobility feasibility).] At every step:
    $e_{v,k} - E^{rem}_{v,k} \ge e^{res}_v$.
  \item[\textbf{A10} (Sequential charge/discharge).] Charging and
    discharging are sequential phases of the station visit.
    An activity plan $\omega$ specifies either a charging session
    ($E^{ch}_\omega \ge 0$, $E^{dis}_\omega = 0$), a discharging (V2G)
    session ($E^{ch}_\omega = 0$, $E^{dis}_\omega \ge 0$), or both
    as consecutive phases.  As shown in
    Proposition~\ref{prop:no-simultaneous} below, any plan with
    simultaneously positive $p^{ch}_{v,k} > 0$ and
    $p^{dis}_{v,k} > 0$ at the same step $k$ is strictly dominated in
    cost under the no-arbitrage price
    condition~\eqref{eq:no-arbitrage-condition}, so no binary mode
    variables are required in the typical operating regime.
  \item[\textbf{A11} (Station as service node).] EVs traverse
    arrival $\to$ queue $\to$ service $\to$ dwell $\to$ egress
    inside the network loading.
  \item[\textbf{A12} (Plug capacity).] $B_{s,k} \le \bar{B}_s$
    at every step.
\end{description}

% ------------------------------------------------------------
\subsection{Microscopic Vehicle Model}
\label{sec:traffic-micro}
% ------------------------------------------------------------

\subsubsection{ICV microscopic state and route choice}

ICV $v$ has state $x^F_{v,k} = (l_{v,k},\, p_v)$ comprising its current
location and chosen path.  For OD $w$ departing at step $k$, the
route-choice problem is
\begin{equation}
  r^*_{v,k} = \arg\min_{r \in \mathcal{R}_w} T_{r,k},
  \label{eq:icv-route-choice}
\end{equation}
where the path travel time is computed by dynamic recursion along
path $r = (a_1, a_2, \ldots, a_m)$.  This equation gives individual
route choice in \emph{fixed-departure mode} where departure step $k$ is
exogenous.  In full DUE, ICVs jointly choose departure step and route to
minimise $\Psi^F_{w,r,k}$ (which includes schedule delay;
see~\S\ref{sec:traffic-psi}).
\begin{equation}
  k_{j+1} = k_j + \left\lceil \frac{T_{a_j, k_j}}{\Delta t} \right\rceil,
  \qquad
  T_{r,k} = \sum_{j=1}^{m} T_{a_j, k_j}.
  \label{eq:path-tt-recursion}
\end{equation}

Under Wardrop time-equilibrium in \emph{fixed-departure mode} the ICV
travel-time condition is
\begin{equation}
  T_{r,k} \begin{cases}
    = T^{\min}_{w,k}   & \text{if } x^F_{w,r,k} > 0, \\
    \ge T^{\min}_{w,k} & \text{if } x^F_{w,r,k} = 0.
  \end{cases}
  \label{eq:icv-wardrop-tt}
\end{equation}
Here $x^F_{w,r,k}=h^F_{w,r,k}$ is the equilibrium route flow at fixed
departure step $k$.  The generalised-cost condition
for \emph{full DUE} (with departure-time endogeneity and schedule delay)
is given by~\eqref{eq:icv-wardrop}--\eqref{eq:icv-wardrop-zero} in
\S\ref{sec:traffic-equilibrium}.

\paragraph{Fuel consumption and emissions.}
Link-level fuel consumption for ICV $v$ on link $a$ is
\begin{equation}
  F_{v,a,k} = L_a \!\left(\alpha_F + \frac{\beta_F}{\bar{v}_{a,k}}\right),
  \label{eq:icv-fuel}
\end{equation}
where $\alpha_F$ (L/km) is the distance-proportional coefficient,
$\beta_F$ (L$\cdot$h/km) the congestion-sensitive coefficient, and
$\bar{v}_{a,k} = L_a / T_{a,k}$ the mean link speed.
Tailpipe CO$_2$ emissions are $G_{v,a,k} = \varepsilon^F F_{v,a,k}$.

\subsubsection{EV microscopic state}

EV $v$ has a four-dimensional state
\begin{equation}
  x^E_{v,k} = \bigl(l_{v,k},\; e_{v,k},\; \sigma_{v,k},\; p_v\bigr),
  \label{eq:ev-state}
\end{equation}
where $l_{v,k}$ is its location, $e_{v,k}$ its battery energy (kWh),
$\sigma_{v,k}$ its behavioural state, and $p_v$ its activity plan.
The behavioural state belongs to
\begin{equation}
  \sigma_{v,k} \in \bigl\{\,
    \textit{driving},\;
    \textit{accessing\_station},\;
    \textit{queueing},\;
    \textit{charging},\;
    \textit{discharging},\;
    \textit{dwelling},\;
    \textit{egressing},\;
    \textit{completed}
  \,\bigr\}.
  \label{eq:ev-states}
\end{equation}

\subsubsection{EV driving energy consumption}

Link-level energy consumption is
\begin{equation}
  E^{drive}_{v,a} = \xi_v\, L_a
  \quad\text{(distance-linear model)},
  \label{eq:ev-energy-linear}
\end{equation}
or, with speed dependence,
\begin{equation}
  E^{drive}_{v,a,k} = g_v\!\bigl(L_a,\, T_{a,k},\, \bar{v}_{a,k}\bigr).
  \label{eq:ev-energy-speed}
\end{equation}
The mean link speed $\bar{v}_{a,k} = L_a / T_{a,k}$ links congestion
directly to energy consumption.

\subsubsection{EV SOC dynamics}

The battery energy evolves as
\begin{equation}
  e_{v,k+1} = e_{v,k}
    - E^{drive}_{v,k}
    + \eta^{ch}_v\, p^{ch}_{v,k}\,\Delta t
    - \frac{p^{dis}_{v,k}\,\Delta t}{\eta^{dis}_v},
  \label{eq:soc-evolution}
\end{equation}
subject to $e^{min}_v \le e_{v,k} \le e^{max}_v$ and the
\emph{mobility feasibility} constraint
\begin{equation}
  e_{v,k} - E^{rem}_{v,k} \ge e^{res}_v,
  \label{eq:soc-feasibility}
\end{equation}
where $E^{rem}_{v,k}$ is the energy required from the current position
to the trip destination and $e^{res}_v$ is the reserve threshold.
Equation~\eqref{eq:soc-feasibility} is the critical coupling between
V2G dispatch and mobility: discharging reduces the feasible plan set
and may necessitate a second charging stop.

\subsubsection{EV charging trigger}

An EV \emph{must} seek a charging stop if the safety constraint~%
\eqref{eq:soc-feasibility} would be violated without one:
\begin{equation}
  e_{v,k} - E^{rem}_{v,k} < e^{res}_v.
  \label{eq:charging-trigger-hard}
\end{equation}
With economic optimisation the trigger extends to any plan that
reduces total cost:
\begin{equation}
  \min_{\omega \in \Omega^{stop}_{v,k}} C^E_{v,\omega,k}
  < C^E_{v,\,\mathrm{no\,stop},\,k},
  \label{eq:charging-trigger-economic}
\end{equation}
meaning an EV with adequate SOC may still enter a station if low energy
prices or V2G revenue make a stop economically attractive.

% ------------------------------------------------------------
\subsection{EV Economic Cost Function}
\label{sec:traffic-ev-cost}
% ------------------------------------------------------------

The total generalised cost of activity plan $\omega$ for OD $w$ departing
at step $k$ is
\begin{equation}
  C^E_{w,\omega,k}
  = C^{travel}_{w,\omega,k}
  + C^{wait}_{w,\omega,k}
  + C^{service}_{w,\omega,k}
  + C^{energy}_{w,\omega,k}
  + C^{deg}_{w,\omega,k}
  + C^{schedule}_{w,\omega,k}
  + \phi^{soc}_{w,\omega,k}.
  \label{eq:ev-gencost}
\end{equation}

\paragraph{Travel-time cost.}
$C^{travel}_{w,\omega,k} = \alpha_t\, T_{\omega,k}$, where $\alpha_t$
(\$/h) is the value of time.

\paragraph{Waiting-time cost.}
If $\omega$ includes station $s$, with $k^{arr}_s$ the arrival step:
\begin{equation}
  C^{wait}_{w,\omega,k} = \alpha_w\, W_{s,k^{arr}_s},
  \label{eq:ev-wait-cost}
\end{equation}
where $W_{s,k}$ is the expected queue waiting time at station $s$ at
step $k$.  This term is zero if $s = \varnothing$.

\paragraph{Service-time cost.}
Charging and discharging are sequential: a vehicle first charges for
duration $T^{ch}_\omega$, then (if V2G) discharges for duration
$T^{dis}_\omega$, occupying a plug for the total time
$T^{service}_\omega = T^{ch}_\omega + T^{dis}_\omega$:
\begin{equation}
  T^{ch}_\omega  = \frac{E^{ch}_\omega}{\eta^{ch}\,\bar{p}^{ch}_s},
  \qquad
  T^{dis}_\omega = \frac{E^{dis}_\omega}{\bar{p}^{dis}_s},
  \label{eq:service-time}
\end{equation}
and $C^{service}_{w,\omega,k} = \alpha_s\, T^{service}_\omega$.

\begin{remark}[Energy-side convention]
  Throughout this section, $p^{ch}_{v,k}$ and $p^{dis}_{v,k}$ are
  \emph{grid-side} power quantities.  The battery receives
  $\eta^{ch}p^{ch}_{v,k}\,\Delta t$ energy per charging step and releases
  $p^{dis}_{v,k}\,\Delta t/\eta^{dis}$ per discharging step, consistent
  with the SOC update~\eqref{eq:soc-evolution}.  The battery-side charged
  energy and grid-side discharged energy for an activity plan are
  \begin{equation}
    E^{ch}_\omega
    = \sum_{\ell\in\mathcal{K}^s_\omega}
      \eta^{ch}\,p^{ch}_{\omega,\ell}\,\Delta t,
    \qquad
    E^{dis}_\omega
    = \sum_{\ell\in\mathcal{K}^s_\omega}
      p^{dis}_{\omega,\ell}\,\Delta t.
    \label{eq:energy-convention}
  \end{equation}
  Charging payments are based on grid-side energy drawn
  ($\sum_\ell p^{ch}_{\omega,\ell}\Delta t = E^{ch}_\omega/\eta^{ch}$);
  V2G revenue is based on grid-side energy delivered ($E^{dis}_\omega$).
  This convention is consistent with the service-time
  formula~\eqref{eq:service-time} and the energy-cost
  term~\eqref{eq:ev-energy-cost}.
\end{remark}
Because discharging extends the total plug-occupancy time
$T^{service}$, it competes directly with queueing vehicles and
enters the cost function without any binary mode variable.

\paragraph{Energy cost.}
Let $\pi^{ch}_{s,\ell}$ and $\pi^{dis}_{s,\ell}$ be the per-kWh
charging price and discharging compensation at station $s$ during
service step $\ell$.  The net energy cost is
\begin{equation}
  C^{energy}_{w,\omega,k}
  = \sum_{\ell \in \mathcal{K}^s_\omega}
    \bigl(\pi^{ch}_{s,\ell}\,p^{ch}_{\omega,\ell}
        - \pi^{dis}_{s,\ell}\,p^{dis}_{\omega,\ell}\bigr)\Delta t,
  \label{eq:ev-energy-cost}
\end{equation}
where charging is a positive cost and V2G revenue is negative.

\paragraph{Battery degradation cost.}
A linear equivalent-throughput model gives
$C^{deg}_{w,\omega,k} = c^{deg}(E^{ch}_\omega + E^{dis}_\omega)$,
where $c^{deg}$ is interpreted as an equivalent cost per kWh of mixed
battery/grid-side throughput.  If degradation is evaluated strictly on
the battery side, replace the discharging term by $E^{dis}_\omega/\eta^{dis}$:
$C^{deg}_{w,\omega,k} = c^{deg,ch}E^{ch}_\omega + c^{deg,dis}E^{dis}_\omega/\eta^{dis}$.
A cycle-depth model $f^{deg}(\mathrm{DoD}_\omega, E^{throughput}_\omega)$
may replace either form for higher fidelity.

\paragraph{Schedule delay cost.}
With desired arrival time $t^*_w$:
\begin{equation}
  C^{schedule}_{w,\omega,k}
  = \gamma_e \max(0,\; t^*_w - t^{arr}_{w,\omega,k})
  + \gamma_l \max(0,\; t^{arr}_{w,\omega,k} - t^*_w).
  \label{eq:schedule-cost}
\end{equation}

\paragraph{SOC infeasibility penalty.}
\begin{equation}
  \phi^{soc}_{w,\omega,k}
  = M \cdot \mathbf{1}\!\left[
      \exists\, \ell : e_{\omega,\ell} - E^{rem}_{\omega,\ell} < e^{res}
    \right],
  \label{eq:soc-penalty}
\end{equation}
where $M$ is a large penalty.  In optimisation formulations the
mobility feasibility~\eqref{eq:soc-feasibility} is imposed as a
hard constraint instead.

\begin{proposition}[Sequential charge/discharge under no-arbitrage prices]
\label{prop:no-simultaneous}
Assume the effective round-trip V2G compensation does not exceed the
charging cost after round-trip efficiency and degradation:
\begin{equation}
  \pi^{dis}_{s,k}\,\eta^{ch}\,\eta^{dis}
  \;\le\;
  \pi^{ch}_{s,k} + c^{deg}.
  \label{eq:no-arbitrage-condition}
\end{equation}
Then for any activity plan $\omega$ with $p^{ch}_{\omega,k} > 0$ and
$p^{dis}_{\omega,k} > 0$ at the same step $k$, there exists a plan
$\omega'$ with strictly lower cost: $C^E_{w,\omega',k} < C^E_{w,\omega,k}$.
\end{proposition}
\begin{proof}
Consider any step $k$ at which $p^{ch}_{\omega,k}>0$ and
$p^{dis}_{\omega,k}>0$.  The net battery-energy increment induced by the
original plan at this step is
\begin{equation}
  \Delta e_k
  =
  \eta^{ch}p^{ch}_{\omega,k}\,\Delta t
  -
  \frac{p^{dis}_{\omega,k}\,\Delta t}{\eta^{dis}}.
  \label{eq:net-soc-increment}
\end{equation}
Construct an alternative plan $\omega'$ that achieves the same net SOC
increment $\Delta e_k$ but avoids simultaneous charging and discharging.
If $\Delta e_k \ge 0$, set
\begin{equation}
  p^{ch}_{\omega',k}
  = \frac{\Delta e_k}{\eta^{ch}\,\Delta t},
  \qquad
  p^{dis}_{\omega',k} = 0.
\end{equation}
If $\Delta e_k < 0$, set
\begin{equation}
  p^{ch}_{\omega',k} = 0,
  \qquad
  p^{dis}_{\omega',k}
  = -\frac{\Delta e_k\,\eta^{dis}}{\Delta t}.
\end{equation}
In both cases one verifies directly that
$\eta^{ch}p^{ch}_{\omega',k}\Delta t - p^{dis}_{\omega',k}\Delta t/\eta^{dis}
= \Delta e_k$, so $\omega'$ preserves the same SOC trajectory at step $k$.

Because the original plan simultaneously charges and discharges, it
contains an internal cycling component that strictly increases total
energy throughput and plug-occupancy time relative to $\omega'$.
Therefore $C^{service}_{\omega'} < C^{service}_{\omega}$ and
$C^{deg}_{\omega'} < C^{deg}_{\omega}$.

Under the no-arbitrage condition~\eqref{eq:no-arbitrage-condition}, the
eliminated cycling component cannot produce a net cost saving after
accounting for round-trip efficiency and degradation, so the energy-cost
term does not increase.  Therefore
$C^E_{w,\omega',k} < C^E_{w,\omega,k}$, proving that no cost-minimising
plan contains simultaneous charging and discharging at the same step.
\end{proof}

\begin{remark}[Price condition]
  Proposition~\ref{prop:no-simultaneous} requires the no-arbitrage
  condition~\eqref{eq:no-arbitrage-condition}: the effective round-trip
  V2G revenue $\pi^{dis}_{s,k}\,\eta^{ch}\eta^{dis}$ must not exceed the
  charging cost $\pi^{ch}_{s,k}$ plus degradation $c^{deg}$.
  When this condition fails (strong V2G premium), simultaneous
  charge/discharge may be cost-improving; binary mode constraints
  $I^{ch}_{v,k}+I^{dis}_{v,k}\le 1$ should then be added.
  Under the no-arbitrage condition, binary variables are not needed and
  the problem remains a pure LP.
\end{remark}

\subsubsection{EV equilibrium condition (fixed-departure special case)}

In fixed-departure mode, the EV Wardrop condition at a given step $k$ is:
\begin{equation}
  h^E_{w,\omega,k} > 0 \;\Rightarrow\;
  C^E_{w,\omega,k} = C^{E,\min}_{w,k}
  := \min_{\omega' \in \Omega_{w,k}} C^E_{w,\omega',k}.
  \label{eq:ev-equilibrium}
\end{equation}
This is the restriction of the full DUE
condition~\eqref{eq:ev-wardrop}--\eqref{eq:ev-wardrop-zero} to a single
departure step; the OD-level minimum $\mu^E_w$ reduces to the step-level
minimum $C^{E,\min}_{w,k}$ when the departure profile is fixed.

% ------------------------------------------------------------
\subsection{Fleet Aggregation Model}
\label{sec:traffic-fleet}
% ------------------------------------------------------------

\subsubsection{ICV departure-rate variables and demand conservation}

The ICV DUE variable is $h^F_{w,r,k} \ge 0$: the number of ICVs of OD
$w$ departing at step $k$ and choosing route $r$.  The
\emph{full-DUE} demand constraint sums across the departure window:
\begin{equation}
  \sum_{k \in \mathcal{K}^{dep}_w}
  \sum_{r \in \mathcal{R}_w}
  h^F_{w,r,k}
  = D^F_w
  \quad \forall\, w.
  \label{eq:icv-conservation-due}
\end{equation}
In \emph{fixed-departure mode} the per-step constraint replaces~\eqref{eq:icv-conservation-due}:
\begin{equation}
  \sum_{r \in \mathcal{R}_w} h^F_{w,r,k} = d^F_{w,k}
  \quad \forall\, w, k.
  \label{eq:icv-conservation}
\end{equation}
The notation $x^F_{w,r,k}$ used in earlier sections (route-flow after
equilibration) equals $h^F_{w,r,k}$ at the fixed point.

\subsubsection{EV activity-plan departure-rate variables and demand conservation}

The EV DUE variable is $h^E_{w,\omega,k} \ge 0$: the number of EVs of
OD $w$ departing at step $k$ and choosing activity plan $\omega$.
The \emph{full-DUE} demand constraint is:
\begin{equation}
  \sum_{k \in \mathcal{K}^{dep}_w}
  \sum_{\omega \in \Omega_{w,k}}
  h^E_{w,\omega,k}
  = D^E_w
  \quad \forall\, w.
  \label{eq:ev-conservation-due}
\end{equation}
In \emph{fixed-departure mode}:
\begin{equation}
  \sum_{\omega \in \Omega_{w,k}} h^E_{w,\omega,k} = d^E_{w,k}
  \quad \forall\, w, k.
  \label{eq:ev-conservation}
\end{equation}

\subsubsection{Fleet-to-link inflow mapping}

Total link inflow $u_{a,k}$ is the sum of contributions from ICV routes,
EV activity plans, and station egress flows:
\begin{equation}
  u_{a,k} = u^F_{a,k} + u^E_{a,k} + u^{egress}_{a,k},
  \label{eq:link-inflow-total}
\end{equation}
where dynamic-incidence indicators $\delta^F_{a,w,r,k}$ and
$\delta^E_{a,w,\omega,k}$ specify whether a given fleet flow enters
link $a$ at step $k$; these are determined endogenously by the
CTM/LTM propagation layer.

\subsubsection{Fleet-to-station arrival and departure}

Station $s$ arrival flow:
\begin{equation}
  u^S_{s,k}
  = \sum_{w} \sum_{\substack{\omega \in \Omega_{w,k} \\ s \in \omega}}
    \delta^{arr}_{s,w,\omega,k}\, x^E_{w,\omega,k},
  \label{eq:station-arrival}
\end{equation}
and station departure flow:
\begin{equation}
  o_{s,k}
  = \sum_{w} \sum_{\substack{\omega \in \Omega_{w,k} \\ s \in \omega}}
    \delta^{dep}_{s,w,\omega,k}\, x^E_{w,\omega,k}.
  \label{eq:station-departure}
\end{equation}
The departure flow $o_{s,k}$ re-enters the road network as
$u^{egress}_{a,k}$ on the station's egress link.

\subsubsection{Class-specific cumulative curves}

Class-specific cumulative entry and exit counts sum to the aggregate:
\begin{equation}
  N^{in}_{a,k} = \sum_{c \in \mathcal{C}} N^{in,c}_{a,k}, \qquad
  N^{out}_{a,k} = \sum_{c \in \mathcal{C}} N^{out,c}_{a,k}.
  \label{eq:class-cumulative}
\end{equation}
Link occupancy at step $k$: $n_{a,k} = N^{in}_{a,k} - N^{out}_{a,k}$.

% ------------------------------------------------------------
\subsection{CTM-Based System-Level Network Loading}
\label{sec:traffic-ctm}
% ------------------------------------------------------------

The \emph{Cell Transmission Model} (CTM, Daganzo 1994) discretises each
link into a sequence of cells, resolving within-link shockwave dynamics and
supporting route-commodity tracking for DUE.
The implementation resides in \texttt{src/ev\_power\_traffic/ctm\_propagation.cpp}.

\subsubsection{Cell discretisation and CFL condition}

Link $a$ is partitioned into $M_a$ cells of equal length
\begin{equation}
  \delta_a = \frac{L_a}{M_a}, \qquad
  M_a = \max\!\left(1,\;\left\lfloor \frac{L_a}{v^{ff}_a\,\Delta t} \right\rfloor\right).
  \label{eq:ctm-cell-length}
\end{equation}
When $L_a \ge v^{ff}_a\,\Delta t$, the choice in~\eqref{eq:ctm-cell-length}
yields $\delta_a = L_a/M_a \ge v^{ff}_a\,\Delta t$, so the
\textbf{CFL condition} $\delta_a \ge v^{ff}_a\,\Delta t$ is satisfied by
construction.  If $L_a < v^{ff}_a\,\Delta t$, no cell partition with
$M_a \ge 1$ can satisfy this condition; the global time step must be
reduced or sub-stepping must be used before applying the CTM.
For the full network:
\begin{equation}
  \Delta t \le \min_{a \in \mathcal{A}^R} \frac{\delta_a}{v^{ff}_a}.
  \label{eq:ctm-global-cfl}
\end{equation}
The jam occupancy per cell is
\begin{equation}
  N^{jam}_{a,m} = k^j_a\,\delta_a
    = \frac{\bar{q}_a(v^{ff}_a + w_a)}{v^{ff}_a\,w_a}\,\delta_a.
  \label{eq:ctm-njam-cell}
\end{equation}

\subsubsection{Cell sending and receiving functions}

\begin{align}
  S_{a,m,k} &= \min\!\left(\frac{v^{ff}_a}{\delta_a}\,n_{a,m,k},\;
                            \bar{q}_a\right)\!\Delta t,
  \label{eq:ctm-send}\\
  R_{a,m,k} &= \min\!\left(\bar{q}_a,\;
                            \frac{w_a}{\delta_a}
                            \bigl(N^{jam}_{a,m} - n_{a,m,k}\bigr)\right)\!\Delta t.
  \label{eq:ctm-recv}
\end{align}

\subsubsection{Inter-cell flow and occupancy update}

The flow from cell $m$ to cell $m+1$:
\begin{equation}
  y_{a,m\to m+1,k} = \min\!\bigl(S_{a,m,k},\; R_{a,m+1,k}\bigr),
  \label{eq:ctm-intercell-flow}
\end{equation}
and the cell occupancy update:
\begin{equation}
  n_{a,m,k+1} = n_{a,m,k}
    + y_{a,m-1\to m,k}
    - y_{a,m\to m+1,k}.
  \label{eq:ctm-cell-update}
\end{equation}

\subsubsection{Route-commodity tracking}

For multi-route networks a per-route fraction is maintained:
\begin{equation}
  \phi^r_{a,m,k} = \frac{n^r_{a,m,k}}{n_{a,m,k} + \varepsilon},
  \qquad n_{a,m,k} = \sum_{r} n^r_{a,m,k},
  \label{eq:ctm-route-fraction}
\end{equation}
and the per-route inter-cell transfer is
$y^r_{a,m\to m+1,k} = \phi^r_{a,m,k}\, y_{a,m\to m+1,k}$.

\begin{remark}[Proportional mixing approximation]
  Equation~\eqref{eq:ctm-route-fraction} propagates route identities by
  proportional mixing: each route receives a share of the inter-cell flow
  proportional to its current fraction of cell occupancy.  This preserves
  aggregate link conservation but does not enforce route-level FIFO
  ordering.  Vehicles from later-departing routes may be assigned to
  earlier exit slots.  Exact route-specific FIFO tracking requires
  per-commodity cumulative curves and is handled by the LTM layer
  (Section~\ref{sec:traffic-ltm}).
\end{remark}

\subsubsection{Inter-link node model}

At node $i$, the inter-link flow from link $a$ to link $b$ is
\begin{equation}
  y_{a\to b,k} = \beta_{a,b,k}
    \min\!\bigl(S^{exit}_{a,k},\; R^{entry}_{b,k}\bigr),
  \label{eq:ctm-node-flow}
\end{equation}
where $S^{exit}_{a,k} = S_{a,M_a-1,k}$,
$R^{entry}_{b,k} = R_{b,0,k}$, and $\beta_{a,b,k}$ is the turning
fraction derived from route fractions or given as exogenous input.

\begin{remark}[Node model limitation]
  Equation~\eqref{eq:ctm-node-flow} is a single-input/single-output
  approximation valid when each junction has at most one upstream link.
  For general merge nodes (multiple incoming links sharing one outgoing
  link), the individual terms $\beta_{a,b,k}\min(S^{exit}_{a,k},
  R^{entry}_{b,k})$ may sum to exceed $R^{entry}_{b,k}$, violating the
  downstream receiving constraint.  A proper node model maximises joint
  throughput subject to sending and receiving constraints simultaneously
  (see equations~\eqref{eq:node-send}--\eqref{eq:node-recv} in the LTM
  section), applying supply-side priority rules at diverge nodes.
\end{remark}
\begin{remark}[Current CTM implementation scope]
  The current \texttt{ctm\_propagation.cpp} implementation applies
  equation~\eqref{eq:ctm-node-flow} with precomputed turning fractions.
  Support for general diverge/merge node solvers (simultaneous
  multi-link capacity allocation) is a planned extension and is not
  yet enabled.
\end{remark}

\subsubsection{CTM-based Dynamic User Equilibrium (DUE)}

The CTM is embedded in a DUE outer loop.  Two solution methods are
available:

\paragraph{MSA (heuristic baseline).}
\begin{align}
  x^{(i+1)} &= x^{(i)} + \frac{1}{i+1}\bigl(\hat{x}^{(i)} - x^{(i)}\bigr),
  \label{eq:msa-update}
\end{align}
where $\hat{x}^{(i)}$ is the all-or-nothing assignment to shortest-time
routes given CTM travel times from iteration $i$.  The relative gap
\begin{equation}
  \epsilon_i = \frac{\sum_r (T^r - T^{\min})\,x^{(i)}_r}
                    {\sum_r T^{\min}\,x^{(i)}_r}
  \label{eq:msa-gap}
\end{equation}
drives convergence.

\paragraph{LP-based system-optimal formulation.}
Section~\ref{sec:traffic-ctm-so-lp} gives the SO-CTM LP that replaces
MSA when system-optimal objectives are acceptable and optimality
certificates are required.

\begin{remark}[CTM vs LTM granularity]
  CTM has $O(M_a |\mathcal{A}| T)$ state variables per simulation; LTM
  has $O(|\mathcal{A}| T)$.  Table~\ref{tab:ctm-ltm-comparison} in
  \S\ref{sec:traffic-ctm-ltm-comparison} gives the full trade-off.
\end{remark}

% ------------------------------------------------------------
\subsection{LTM-Based Meso-Micro Network Loading}
\label{sec:traffic-ltm}
% ------------------------------------------------------------

The \emph{Link Transmission Model} (LTM, Yperman 2007) propagates traffic
using cumulative entry and exit curves rather than per-cell densities.
It is the primary propagation layer for EV meso-micro coupling because it
provides exact FIFO trajectory reconstruction at $O(|\mathcal{A}|T)$ cost.

\subsubsection{Cumulative curves and occupancy}

For link $a$, define cumulative entry and exit counts:
\begin{equation}
  u_{a,k} = N^{in}_{a,k+1} - N^{in}_{a,k}, \qquad
  v_{a,k} = N^{out}_{a,k+1} - N^{out}_{a,k}.
  \label{eq:ltm-flows}
\end{equation}
Link occupancy: $n_{a,k} = N^{in}_{a,k} - N^{out}_{a,k}$.

\subsubsection{Free-flow and backward-wave delays}

\begin{equation}
  \tau^{ff}_a = \left\lceil \frac{L_a}{v^{ff}_a \Delta t} \right\rceil,
  \qquad
  \tau^{bw}_a = \left\lceil \frac{L_a}{w_a \Delta t} \right\rceil.
  \label{eq:tau-ff-bw}
\end{equation}

\begin{remark}[Initial history curves]
  The sending and receiving functions~\eqref{eq:ltm-sending}--\eqref{eq:ltm-receiving}
  reference $N^{in}_{a,k-\tau^{ff}_a}$ and $N^{out}_{a,k-\tau^{bw}_a}$
  for $k < \tau$.  For an empty initial network:
  $N^{in}_{a,k} = N^{out}_{a,k} = 0$ for all $k \le 0$.
  Non-zero pre-history (e.g., background demand already en route)
  must be supplied externally before the first simulation step.
\end{remark}

\subsubsection{LTM sending and receiving functions}

\begin{align}
  S_{a,k} &= \min\!\left(
    N^{in}_{a,\,k-\tau^{ff}_a} - N^{out}_{a,k},\;
    \bar{q}_a\,\Delta t
  \right),
  \label{eq:ltm-sending}\\
  R_{a,k} &= \min\!\left(
    N^{out}_{a,\,k-\tau^{bw}_a} + N^{jam}_a - N^{in}_{a,k},\;
    \bar{q}_a\,\Delta t
  \right).
  \label{eq:ltm-receiving}
\end{align}
The boundary flow constraints are
$0 \le v_{a,k} \le S_{a,k}$ and $0 \le u_{a,k} \le R_{a,k}$.

\begin{remark}[Discrete-offset convention]
  In the current implementation, entry flows at step $k$ are written to
  $N^{in}_{a,k+1}$ (flows are increments between steps), and exit flows
  are written after the sending calculation.  Depending on the update
  ordering, the first observable exit index in the code may appear at
  $k + \tau^{ff}_a + 1$ or $k + \tau^{ff}_a + 2$, rather than the
  theoretical $k + \tau^{ff}_a$.  This is an implementation-indexing
  convention arising from the entry-increment and outflow write-back
  ordering, not an intrinsic LTM property.  The theoretical minimum
  travel time remains $\tau^{ff}_a\,\Delta t$.
\end{remark}

\subsubsection{LTM node model}

For intersection node $i$ with incoming link set $\mathcal{A}^{in}(i)$
and outgoing link set $\mathcal{A}^{out}(i)$:
\begin{align}
  \sum_{b \in \mathcal{A}^{out}(i)} y_{a,b,k} &\le S_{a,k},
    \quad a \in \mathcal{A}^{in}(i),
    \label{eq:node-send}\\
  \sum_{a \in \mathcal{A}^{in}(i)} y_{a,b,k} &\le R_{b,k},
    \quad b \in \mathcal{A}^{out}(i).
    \label{eq:node-recv}
\end{align}
Link flows satisfy
\begin{align}
  v_{a,k} &= \sum_{b \in \mathcal{A}^{out}(i)} y_{a,b,k}, \\
  u_{b,k} &= \sum_{a \in \mathcal{A}^{in}(i)} y_{a,b,k}
              + d^{orig}_{b,k}
              + o_{s\to b,k},
  \label{eq:node-entry}
\end{align}
where $d^{orig}_{b,k}$ is exogenous origin demand and
$o_{s\to b,k}$ is the egress flow entering link $b$ from station $s$.

\subsubsection{FIFO-consistent vehicle trajectory reconstruction}
\label{sec:traffic-fifo-reconstruction}

The LTM layer is mesoscopic and does not explicitly store individual
trajectories.  They are recovered by FIFO-consistent cumulative-count
inversion.

For a vehicle entering link $a$ during step $k$, its FIFO rank lies in
the cumulative interval
\begin{equation}
  \xi_{v,a} \in \bigl(N^{in}_{a,k},\; N^{in}_{a,k+1}\bigr].
\end{equation}
If vehicles are represented individually, the $j$-th vehicle entering
during step $k$ is assigned $\xi_{v,a} = N^{in}_{a,k} + j$; for
continuous flows a representative rank within the interval is used.
The exit step is then
\begin{equation}
  k^{out}_{v,a}
  = \min\!\bigl\{k : N^{out}_{a,k} \ge \xi_{v,a}\bigr\}.
  \label{eq:fifo-exit}
\end{equation}
Repeating over the activity path yields the full trajectory:
$\mathcal{T}_v = \{(a,\, k^{in}_{v,a},\, k^{out}_{v,a}) : a \in p_v\}$.

Station ranks are assigned analogously in cumulative intervals,
where $U_{s,k}$ is cumulative arrivals, $B^{cum}_{s,k}$ cumulative service
starts, $C^{cum}_{s,k}$ cumulative completions, and $O_{s,k}$ cumulative
departures at station $s$ by step $k$:
arrivals during step $k$ occupy $(U_{s,k},\, U_{s,k+1}]$,
service starts occupy $(B^{cum}_{s,k},\, B^{cum}_{s,k+1}]$, and
completions occupy $(C^{cum}_{s,k},\, C^{cum}_{s,k+1}]$.
For the vehicle arriving as the $\xi^{arr}_{v,s}$-th at station $s$:
\begin{align}
  k^{start}_{v,s} &= \min\!\bigl\{k : B^{cum}_{s,k} \ge \xi^{arr}_{v,s}\bigr\},
    \label{eq:charge-start}\\
  k^{end}_{v,s}   &= \min\!\bigl\{k : C^{cum}_{s,k} \ge \xi^{start}_{v,s}\bigr\},
    \label{eq:charge-end}\\
  k^{dep}_{v,s}   &= \min\!\bigl\{k : O_{s,k} \ge \xi^{end}_{v,s}\bigr\}.
    \label{eq:station-dep}
\end{align}
This links aggregate LTM states to individual EV records without a full
microscopic simulation.

% ------------------------------------------------------------
\subsection{Charging Stations as Endogenous Service Nodes}
\label{sec:traffic-charging-station}
% ------------------------------------------------------------

\subsubsection{Station state and dynamics}

Station $s$ tracks three aggregate populations:
$Q_{s,k}$ (queueing), $B_{s,k}$ (in service), $H_{s,k}$ (post-service dwell).
The state transitions are
\begin{align}
  Q_{s,k+1} &= Q_{s,k} + u^S_{s,k} - b_{s,k},
    \label{eq:st-Q}\\
  B_{s,k+1} &= B_{s,k} + b_{s,k}   - c_{s,k},
    \label{eq:st-B}\\
  H_{s,k+1} &= H_{s,k} + c_{s,k}   - o_{s,k}.
    \label{eq:st-H}
\end{align}

\subsubsection{Service admission}

\begin{equation}
  b_{s,k} \le Q_{s,k} + u^S_{s,k},
  \qquad
  B_{s,k} + b_{s,k} \le \bar{B}_s.
  \label{eq:admission-constraints}
\end{equation}

\subsubsection{Service completion}

\begin{equation}
  c_{s,k} \le B_{s,k}, \qquad c_{s,k} \le \mu^{step}_{s,k},
  \label{eq:completion-constraints}
\end{equation}
where in the fluid station model
\begin{equation}
  \mu^{step}_{s,k}
  = \frac{\bar{B}_s\,\Delta t}{\bar{T}^{service}_{s,k}}
  \label{eq:mu-step}
\end{equation}
is the per-step service completion capacity (veh/step).  In the
individual-vehicle mode, $c_{s,k} = \sum_v \mathbf{1}[k^{end}_{v,s}=k]$.

\subsubsection{Station departure}

\begin{equation}
  o_{s,k} \le H_{s,k} + c_{s,k},
  \qquad
  o_{s,k} \le R^{egress}_{s,k},
  \label{eq:departure-constraints}
\end{equation}
where $R^{egress}_{s,k}$ is the receiving capacity of the downstream
road link.  An overcrowded downstream link prevents EV departure even
after service completion.

\subsubsection{Queue spillback to road}

If $\bar{Q}_s < \infty$ is the queue capacity, the access-link receiving
capacity is reduced as the queue fills:
\begin{equation}
  R^{access}_{s,k}
  = R^0_{s,k}\,\left(1 - \frac{Q_{s,k}}{\bar{Q}_s}\right)_{\!+},
  \label{eq:spillback}
\end{equation}
so that $R^{access}_{s,k} = 0$ when $Q_{s,k} \ge \bar{Q}_s$.
This feedback propagates station congestion back into the road network.

\begin{remark}[Spillback and egress feedback: partially wired]
  Equations~\eqref{eq:departure-constraints} and~\eqref{eq:spillback}
  specify the egress-capacity bound $R^{egress}_{s,k}$ and the
  access-link capacity reduction, respectively. In the current
  implementation, CTM applies a cumulative-arrival proxy, and LTM uses
  queue-based receiving reduction. Thus, station arrivals do feed back
  into upstream congestion. However, this is still a partial formulation
  relative to the full bidirectional theoretical model.
\end{remark}

% ------------------------------------------------------------
\subsection{Charging and V2G Feedback on Traffic}
\label{sec:traffic-feedback}
% ------------------------------------------------------------

Charging and V2G decisions affect traffic through three channels.

\paragraph{Channel 1: Route and plan choice.}
Low energy prices or high V2G compensation may induce EVs to choose
a longer route that passes a favoured station:
\begin{equation}
  C^E_{\omega_B,k} < C^E_{\omega_A,k}
  \quad\Rightarrow\quad
  x^E_{w,\omega_B,k} > 0,
  \label{eq:route-switching}
\end{equation}
changing link inflows $u_{a,k}$ and thus congestion.

\paragraph{Channel 2: SOC depletion and secondary charging.}
V2G discharge reduces $e_{v,k}$, potentially triggering the
safety condition~\eqref{eq:charging-trigger-hard} later in the trip.
A secondary charging stop adds an extra station visit, increases
$u^S_{s',k}$ at a downstream station $s'$, and extends total trip time.

\paragraph{Channel 3: Dwell time and egress.}
Charging duration $T^{ch}_\omega = E^{ch}_\omega / (\eta^{ch}\bar{p}^{ch}_s)$
directly determines how long a vehicle occupies a plug.  Larger charging
energy increases $B_{s,k}$, reduces $\bar{B}_s - B_{s,k}$, and can grow
the queue $Q_{s,k}$, which eventually activates the spillback
mechanism~\eqref{eq:spillback}.

\paragraph{Summary of coupling.}
\begin{equation}
  \underbrace{(E^{ch},\, E^{dis},\, \pi^{ch},\, \pi^{dis})}_{\text{charging/V2G strategy}}
  \;\;\longrightarrow\;\;
  \underbrace{(u^S_{s,k},\; o_{s,k},\; u_{a,k})}_{\text{station and link flows}}
  \;\;\longrightarrow\;\;
  \underbrace{(N^{in},\, N^{out},\, T_{a,k})}_{\text{LTM/CTM state}}
  \;\;\longrightarrow\;\;
  \underbrace{C^E_{w,\omega,k}}_{\text{EV cost feedback}}.
  \label{eq:feedback-chain}
\end{equation}

% ------------------------------------------------------------
\subsection{Dynamic Equilibrium Model}
\label{sec:traffic-equilibrium}
% ------------------------------------------------------------

This module supports two assignment modes:
(i)~\textbf{Full DUE}: departure time, route, and EV activity plan are
endogenous choices under OD total-demand constraints~\eqref{eq:icv-conservation-due}--\eqref{eq:ev-conservation-due};
equilibrium costs are compared across all departure steps and alternatives.
(ii)~\textbf{Fixed-departure DUE}: the departure profile $d_{w,k}$ is
given; only route/activity-plan choice is equilibrated within each
departure step under constraints~\eqref{eq:icv-conservation}--\eqref{eq:ev-conservation}.

\begin{remark}[Implementation scope: Full DUE for EV windows]
  The CTM-DUE solver supports both modes.  With
  \texttt{DUEOptions::due\_mode = DUEMode::FullEndogenous}, each EV demand
  with non-empty \texttt{EVDemand::departure\_window\_steps} is assigned over
  the full $(k,\omega)$ window: the all-or-nothing direction minimises
  generalized cost across departure step, route/activity plan, charging cost,
  V2G credit, and schedule-delay penalties.  The result records both route
  totals and disaggregated
  \texttt{flow\_by\_demand\_departure\_route}.  Fixed-departure mode remains
  the default for backward compatibility.  Multi-class ICV departure-time
  endogeneity is still limited to fixed \texttt{ICVDemand::departure\_step};
  extending the ICV multi-class branch to cross-step departure windows is a
  remaining implementation extension.
\end{remark}

Unless stated otherwise, the mathematical formulations in this section
describe full DUE~(i); benchmark cases explicitly state whether they use
full-endogenous or fixed-departure mode.

\subsubsection{Dynamic Network Loading operator}

The DNL operator maps departure-rate flows to network states:
\begin{equation}
  \mathcal{D} : (h^F, h^E)
  \;\mapsto\;
  (T,\, W,\, Q,\, o,\, N^{in},\, N^{out}).
  \label{eq:dnl-operator}
\end{equation}

\subsubsection{Experienced-cost operators}
\label{sec:traffic-psi}

After DNL produces network states, the \emph{experienced-cost operators}
convert states to generalised path costs:
\begin{align}
  \Psi^F_{w,r,k}(h) &:=
    \alpha_t\,T^{exp}_{w,r,k}(h) + C^{schedule,F}_{w,r,k}(h),
    \label{eq:psi-icv} \\
  \Psi^E_{w,\omega,k}(h) &:=
    C^E_{w,\omega,k}(h),
    \label{eq:psi-ev}
\end{align}
where $T^{exp}_{w,r,k}(h)$ is the FIFO-reconstructed experienced travel
time of a departure-$k$ vehicle on route $r$,
$C^E_{w,\omega,k}$ is the EV generalised cost from~\eqref{eq:ev-gen-cost},
and the ICV schedule-delay penalty is
\begin{equation}
  C^{schedule,F}_{w,r,k}(h)
  = \gamma_e \max\!\bigl(0,\, t^*_w - t^{arr}_{w,r,k}\bigr)
  + \gamma_l \max\!\bigl(0,\, t^{arr}_{w,r,k} - t^*_w\bigr),
  \label{eq:icv-schedule-delay}
\end{equation}
with $t^{arr}_{w,r,k} = k\Delta t + T^{exp}_{w,r,k}$, desired arrival
time $t^*_w$, and early/late penalties $\gamma_e, \gamma_l \ge 0$.
Setting $\gamma_e = \gamma_l = 0$ recovers the travel-time-only special
case.

\subsubsection{ICV and EV choice operators}

The DUE minimum costs are OD-level minima across \emph{all} departure
steps and alternatives in the departure window:
\begin{equation}
  \mu^F_w := \min_{k \in \mathcal{K}^{dep}_w,\; r \in \mathcal{R}_w}
    \Psi^F_{w,r,k}(h),
  \label{eq:icv-min-cost}
\end{equation}
\begin{equation}
  \mu^E_w := \min_{k \in \mathcal{K}^{dep}_w,\; \omega \in \Omega_{w,k}}
    \Psi^E_{w,\omega,k}(h).
  \label{eq:ev-min-cost}
\end{equation}
The ICV assignment operator $\mathcal{A}^F$ and EV assignment operator
$\mathcal{A}^E$ map experienced costs to departure-rate flows.
Under all-or-nothing (AON) assignment for full DUE,
the entire OD demand is assigned to the minimiser across all $(k, r)$
(resp.\ $(k, \omega)$):
\begin{equation}
  \hat{h}^F_{w,r,k} =
  \begin{cases}
    D^F_w & (r,k) \in \arg\min_{r',k'} \Psi^F_{w,r',k'}, \\
    0     & \text{otherwise},
  \end{cases}
  \label{eq:aon-icv}
\end{equation}
\begin{equation}
  \hat{h}^E_{w,\omega,k} =
  \begin{cases}
    D^E_w & (\omega,k) \in \arg\min_{\omega',k'} \Psi^E_{w,\omega',k'}, \\
    0     & \text{otherwise}.
  \end{cases}
  \label{eq:aon-ev}
\end{equation}
Ties are broken by uniform splitting.  A logit stochastic model gives
$P^E_{w,\omega,k} \propto \exp(-\theta\,\Psi^E_{w,\omega,k})$ and is
normalised so that $\sum_{k,\omega} h^E_{w,\omega,k} = D^E_w$.

In \emph{fixed-departure mode}, the minimisation in~\eqref{eq:icv-min-cost}--\eqref{eq:aon-ev} is restricted to the given
departure step $k$, and $D^F_w, D^E_w$ are replaced by $d^F_{w,k}, d^E_{w,k}$.

\subsubsection{Wardrop DUE conditions}

At a full DUE $h^*$, no vehicle can reduce its experienced cost by
unilaterally changing departure time, route, or activity plan:
\begin{align}
  h^{F,*}_{w,r,k} > 0 &\;\Rightarrow\;
    \Psi^F_{w,r,k}(h^*) = \mu^{F,*}_w,
    \label{eq:icv-wardrop} \\
  h^{F,*}_{w,r,k} = 0 &\;\Rightarrow\;
    \Psi^F_{w,r,k}(h^*) \ge \mu^{F,*}_w,
    \label{eq:icv-wardrop-zero} \\
  h^{E,*}_{w,\omega,k} > 0 &\;\Rightarrow\;
    \Psi^E_{w,\omega,k}(h^*) = \mu^{E,*}_w,
    \label{eq:ev-wardrop} \\
  h^{E,*}_{w,\omega,k} = 0 &\;\Rightarrow\;
    \Psi^E_{w,\omega,k}(h^*) \ge \mu^{E,*}_w.
    \label{eq:ev-wardrop-zero}
\end{align}
Here $\mu^{F,*}_w$ and $\mu^{E,*}_w$ are the equilibrium (minimised)
OD costs from~\eqref{eq:icv-min-cost}--\eqref{eq:ev-min-cost} evaluated
at $h^*$.

\begin{remark}[Relationship to fixed-departure condition]
  The condition in~\eqref{eq:ev-equilibrium} is the fixed-departure
  special case of~\eqref{eq:ev-wardrop}--\eqref{eq:ev-wardrop-zero}
  with the minimum taken only over $\omega$ at a given $k$.
\end{remark}

\subsubsection{Fixed-point equilibrium}

The coupled system is at equilibrium when
\begin{equation}
  (h^{F,*},\, h^{E,*})
  = \mathcal{A}\!\bigl(\mathcal{D}(h^{F,*},\, h^{E,*})\bigr),
  \label{eq:fixed-point}
\end{equation}
i.e.\ given the congestion and queues produced by the current departure
rates, no vehicle has an incentive to change its departure time, route,
or activity plan.

% ------------------------------------------------------------
\subsection{Solution Algorithm: MSA for Mixed Fleet DUE}
\label{sec:traffic-msa}
% ------------------------------------------------------------

Starting from initial flows $h^{F,(0)}, h^{E,(0)}$ satisfying the
demand constraints, each iteration $i$ proceeds as:

\begin{enumerate}
  \item \textbf{DNL forward pass:}
    $(T^{(i)}, W^{(i)}, Q^{(i)}, o^{(i)})
    = \mathcal{D}(h^{F,(i)}, h^{E,(i)})$.
  \item \textbf{Compute experienced costs:}
    Evaluate $\Psi^{F,(i)}_{w,r,k}$ and $\Psi^{E,(i)}_{w,\omega,k}$
    via~\eqref{eq:psi-icv}--\eqref{eq:psi-ev} for all $(w,r,k)$ and
    $(w,\omega,k)$.
  \item \textbf{AON assignment (cross-$k$):}
    For full DUE, assign each OD's total demand to the single minimiser
    across all departure steps and alternatives:
    \begin{equation}
      \hat{h}^{F,(i)}_{w,r,k}
      = \begin{cases}
          D^F_w & (r,k) \in \arg\min_{r',k'} \Psi^{F,(i)}_{w,r',k'}, \\
          0     & \text{otherwise},
        \end{cases}
      \label{eq:msa-aon-icv}
    \end{equation}
    and analogously for $\hat{h}^{E,(i)}_{w,\omega,k}$ with $D^E_w$.
    In fixed-departure mode, the minimisation is restricted to each $k$
    with demand $d^F_{w,k}$ (resp.\ $d^E_{w,k}$).
  \item \textbf{MSA update} (step size $\lambda_i = 1/(i+1)$):
    \begin{equation}
      h^{F,(i+1)} = h^{F,(i)} + \lambda_i\bigl(\hat{h}^{F,(i)} - h^{F,(i)}\bigr),
      \quad
      h^{E,(i+1)} = h^{E,(i)} + \lambda_i\bigl(\hat{h}^{E,(i)} - h^{E,(i)}\bigr).
      \label{eq:msa-mixed-update}
    \end{equation}
  \item \textbf{Convergence check.}
    The \emph{ICV relative gap} (OD-level minimum):
    \begin{equation}
      \epsilon^F_i
      = \frac{
          \displaystyle\sum_{w,k,r}
            h^{F,(i)}_{w,r,k}\,
            \bigl(\Psi^{F,(i)}_{w,r,k} - \mu^{F,(i)}_w\bigr)
        }{
          \displaystyle\sum_w D^F_w\,\mu^{F,(i)}_w
        },
      \label{eq:icv-gap}
    \end{equation}
    where $\mu^{F,(i)}_w = \min_{k,r} \Psi^{F,(i)}_{w,r,k}$.
    The \emph{EV relative gap}:
    \begin{equation}
      \epsilon^E_i
      = \frac{
          \displaystyle\sum_{w,k,\omega}
            h^{E,(i)}_{w,\omega,k}\,
            \bigl(\Psi^{E,(i)}_{w,\omega,k} - \mu^{E,(i)}_w\bigr)
        }{
          \displaystyle\sum_w D^E_w\,\bigl|\mu^{E,(i)}_w\bigr| + \varepsilon
        },
      \label{eq:ev-gap}
    \end{equation}
    where $\mu^{E,(i)}_w = \min_{k,\omega} \Psi^{E,(i)}_{w,\omega,k}$
    and $\varepsilon > 0$ guards against a near-zero denominator when
    EV costs include V2G revenue that can make $\mu^{E,(i)}_w$ negative.
    Stop when $\max(\epsilon^F_i, \epsilon^E_i) < \epsilon^{tol}$.
\end{enumerate}

% ------------------------------------------------------------
\subsection{Key Output Metrics}
\label{sec:traffic-outputs}
% ------------------------------------------------------------

\paragraph{Traffic metrics.}
Total travel time in network: $\mathrm{TTT} = \sum_{a,k} n_{a,k}\,\Delta t$.
Total system travel time:
$\mathrm{TSTT} = \sum_{w,r,k} x^F_{w,r,k} T_{w,r,k}
               + \sum_{w,\omega,k} x^E_{w,\omega,k} T_{w,\omega,k}$.
Network throughput: $\sum_{a \in \mathcal{A}^{sink}} N^{out}_{a,T}$.

\paragraph{Station metrics.}
Maximum queue $Q^{max}_s = \max_k Q_{s,k}$,
average waiting time $\bar{W}_s = \sum_k Q_{s,k}\,\Delta t / \sum_k u^S_{s,k}$,
and plug utilisation $\rho_s = \sum_k B_{s,k} / (\bar{B}_s T)$.
The step-level expected waiting time used in the EV cost~\eqref{eq:ev-wait-cost}
is estimated by the fluid approximation
$W_{s,k} \approx Q_{s,k} / \mu^{hour}_{s,k}$,
where
$\mu^{hour}_{s,k} = \bar{B}_s / \bar{T}^{service}_{s,k}$
(measured in veh/hour when $\bar{T}^{service}_{s,k}$ is in hours)
gives $W_{s,k}$ in hours.  The equivalent per-step capacity is
$\mu^{step}_{s,k} = \mu^{hour}_{s,k}\,\Delta t = \bar{B}_s\Delta t / \bar{T}^{service}_{s,k}$
(veh/step).

\paragraph{EV energy metrics.}
Total charging energy $E^{ch}_{tot} = \sum_{s,k} P^{ch}_{s,k}\,\Delta t$,
total V2G energy $E^{dis}_{tot} = \sum_{s,k} P^{dis}_{s,k}\,\Delta t$,
net load energy $E^{net}_{tot} = E^{ch}_{tot} - E^{dis}_{tot}$,
and V2G potential:
\begin{equation}
  P^{v2g,pot}_{s,k} =
  \sum_{v \in \mathcal{V}^{dwell}_{s,k}}
  \min\!\left(\bar{p}^{dis}_v,\;
    \frac{e_{v,k} - e^{res}_v}{\Delta t}\,\eta^{dis}_v\right).
  \label{eq:v2g-potential}
\end{equation}

Equation~\eqref{eq:v2g-potential} gives the \emph{one-step} V2G power
potential for vehicles currently dwelling at the station.  Multi-period
V2G dispatch additionally requires intertemporal constraints on
SOC~\eqref{eq:soc-evolution} and remaining dwell time
$k^{dep}_{v,s} - k$ to ensure feasibility across the full dwell window.

\paragraph{Station net load.}
Station charging and discharging powers are
$P^{ch}_{s,k} = \sum_v p^{ch}_{v,k}$ and $P^{dis}_{s,k} = \sum_v p^{dis}_{v,k}$,
giving net load $P^{net}_{s,k} = P^{ch}_{s,k} - P^{dis}_{s,k}$ and
bus-level EV injection
$P^{EV}_{b,k} = \sum_{s:\mathrm{bus}(s)=b} P^{net}_{s,k}$.

% ------------------------------------------------------------
\subsection{Correctness Checklist}
\label{sec:traffic-correctness}
% ------------------------------------------------------------

The following invariants must hold throughout every simulation:

\begin{enumerate}
  \item \textbf{Vehicle conservation.} $N^{out}_{a,k} \le N^{in}_{a,k}$
        for all $a, k$.
  \item \textbf{Road capacity.} Admitted link inflows and outflows satisfy
        $u_{a,k} \le \bar{q}_a\,\Delta t$ and
        $v_{a,k} \le \bar{q}_a\,\Delta t$.  Exogenous origin demand may
        exceed this value; any excess is stored in an origin queue and
        admitted in subsequent steps.
  \item \textbf{Jam storage.} $n_{a,k} \le N^{jam}_a$.
  \item \textbf{Plug capacity.} $B_{s,k} \le \bar{B}_s$.
  \item \textbf{SOC feasibility.}
        $e^{min}_v \le e_{v,k} \le e^{max}_v$ and
        $e_{v,k} - E^{rem}_{v,k} \ge e^{res}_v$.
  \item \textbf{FIFO.} If $k^{in}_{i,a} \le k^{in}_{j,a}$ then
        $k^{out}_{i,a} \le k^{out}_{j,a}$.
  \item \textbf{Sequential charge/discharge}
        (Proposition~\ref{prop:no-simultaneous}).
        $p^{ch}_{v,k}\cdot p^{dis}_{v,k} = 0$ at every step $k$
        when the no-arbitrage condition~\eqref{eq:no-arbitrage-condition}
        holds; sequential behaviour then follows from cost optimality
        alone, without binary mode variables.  When V2G compensation
        substantially exceeds charging cost after round-trip efficiency,
        binary mode constraints $I^{ch}_{v,k}+I^{dis}_{v,k}\le 1$ should
        be added.
\end{enumerate}

% ------------------------------------------------------------
\subsection{Numerical Results}
\label{sec:traffic-numerical-results}
% ------------------------------------------------------------

The following results are produced by five test suites, all assertions
passing:
\begin{itemize}
  \item \texttt{test\_ltm\_stress.cpp} — 20 cases, 1532 assertions
        (LTM wave propagation, CFL, conservation).
  \item \texttt{test\_ctm\_ltm\_comparison.cpp} — 5 cases, 118 assertions
        (CTM vs.\ LTM throughput and travel-time accuracy).
  \item \texttt{test\_ev\_power\_traffic\_report.cpp} — 4 cases, 11 assertions
        (fleet assignment and V2G simulation with per-step output).
  \item \texttt{test\_ev\_power\_traffic\_joint\_opt\_d.cpp} — 10 cases, 93 assertions
        (joint social-welfare optimizer, FormD).
  \item \texttt{test\_ev\_power\_traffic\_largescale.cpp} — 3 cases, 14 assertions
        (10,000-EV scalability benchmark on IEEE-118 power grid).
\end{itemize}

Results are organised in two groups.  The first group
(\S\S\ref{sec:traffic-nr-single-link}--\ref{sec:traffic-nr-ltm-ctm-comparison})
contains \textbf{implementation diagnostics}: tests that characterise
known numerical artefacts (integer-delay, CFL) and compare CTM vs.\ LTM
behaviour.  These are diagnostic results for implementation quality, not
validation of the DUE equilibrium model.  The second group
(\S\S\ref{sec:traffic-nr-ev-penetration}--\ref{sec:traffic-nr-largescale})
contains \textbf{model validation tests}: fleet assignment, V2G
optimality, and scalability results that exercise the DUE and
joint-optimisation layers.

\subsubsection{Implementation diagnostics}

These tests characterise the discrete-time LTM and CTM numerical
behaviour; they are not model-validation results for the DUE equilibrium.
Production DUE runs should use time steps and link discretisations for
which the artefacts quantified below (sustained-capacity degradation,
travel-time overestimation at low CFL numbers) are controlled or verified
by time-step refinement tests.

% ··············································
\subsubsection{Single-link LTM: propagation delay and CFL}
\label{sec:traffic-nr-single-link}
% ··············································

For a single urban arterial link ($L = 2\,\text{km}$, $v^{ff} = 50\,\text{km/h}$,
$w = 16.7\,\text{km/h}$, $\bar{q} = 900\,\text{veh/h}$, $\Delta t = 2\,\text{min}$):
\begin{equation*}
  v^{ff}\,\Delta t = 50 \times \tfrac{1}{30} = 1.67\,\text{km} < L = 2\,\text{km}
  \quad\text{[CFL satisfied]},
\end{equation*}
\begin{equation*}
  \tau^{ff} = \left\lceil\tfrac{2}{1.67}\right\rceil = 2, \quad
  \tau^{bw} = \left\lceil\tfrac{2}{0.556}\right\rceil = 4, \quad
  N^{jam} \approx 144\,\text{veh}, \quad
  \text{cap\_step} = 30\,\text{veh}.
\end{equation*}
A 20-vehicle pulse (well below $\text{cap\_step} = 30$) enters at step 0.
The LTM (Group XXI) shows:
\begin{itemize}
  \item $N^{in}[T] = 20$ (full pulse absorbed),
  \item $N^{out}[T] = 20$ (all vehicles drain; $T = 30 \gg \tau^{ff}+\tau^{bw}=6$),
  \item cumulative curves are monotone non-decreasing for all $k$,
  \item CTM gives identical conservation ($N^{in} = N^{out} = 20$).
\end{itemize}

% ··············································
\subsubsection{Freeway sustained-capacity: LTM vs CTM throughput}
% ··············································

Freeway link ($L = 10\,\text{km}$, $v^{ff} = 100\,\text{km/h}$,
$w = 25\,\text{km/h}$, $\bar{q} = 2000\,\text{veh/h}$, $\Delta t = 5\,\text{min}$):
\begin{equation*}
  v^{ff}\,\Delta t = 8.33\,\text{km} < L = 10\,\text{km}
  \quad\text{[CFL satisfied]},
\end{equation*}
\begin{equation*}
  \tau^{ff} = 2, \quad \tau^{bw} = 5, \quad
  N^{jam} = 1000\,\text{veh}, \quad
  \text{cap\_step} = 166.7\,\text{veh}.
\end{equation*}
With $T = 80$ steps of sustained capacity demand (Group XXII):
\begin{center}
\begin{tabular}{lcc}
\toprule
Model & Exit vehicles & Fraction of $T\cdot\text{cap\_step}$ \\
\midrule
CTM & 13\,300 & 99.7\% \\
LTM & 8\,833  & 66.2\% \\
\bottomrule
\end{tabular}
\end{center}
The CTM nearly matches $T \cdot \text{cap\_step} = 13\,333$; in the current
single-cell implementation, the inject-propagate update ordering allows a
vehicle entering at step $k$ to exit at the same step (a zero-delay
implementation artefact that occurs when too few cells are used).
The LTM throughput is lower due to a discretisation-induced
storage--delay mismatch: with $\tau^{ff}=2$ and $\tau^{bw}=5$, the
discrete delay product gives
$\bar{q}\,\Delta t\,(\tau^{ff}+\tau^{bw}) = 166.7 \times 7 = 1167\,\text{veh}$,
which exceeds $N^{jam} = 1000\,\text{veh}$.  This causes an oscillatory
fill/drain pattern.  A crude fill/block-cycle approximation gives
throughput on the order of $6/(6+5)\approx 55\%$ of capacity; the
finite-horizon implementation reports 66.2\% because start-up transients,
terminal drainage, and the discrete update convention allow additional
vehicles to exit before the horizon closes.  The key point is not the
exact fraction but the presence of a storage--delay mismatch under the
coarse $\Delta t=5\,\text{min}$ discretisation.
\textbf{This is a numerical artefact} that should be controlled by
reducing $\Delta t$, using the corrected discrete jam density
$N^{jam,disc}_a = \bar{q}_a\,\Delta t\,(\tau^{ff}_a+\tau^{bw}_a)$,
or a corrected update convention --- it is not an intrinsic LTM property.

% ··············································
\subsubsection{CFL-violation pathology}
% ··············································

Intentionally violating CFL by using $L = 0.5\,\text{km}$ with the
same freeway parameters ($v^{ff}\,\Delta t = 8.33\,\text{km} \gg L$,
Group XXIII):
\begin{center}
\begin{tabular}{lcc}
\toprule
Model & Average throughput & Fraction of cap \\
\midrule
LTM & 150 veh/h & 7.5\% (collapses) \\
CTM coarse (same $\Delta t$) & 2000 veh/h & 100\% (physically invalid) \\
CTM fine ($\Delta t_{ctm}=0.2\,\text{min}$, CFL OK) & 1999 veh/h & 100\% (correct) \\
\bottomrule
\end{tabular}
\end{center}
The LTM collapses because $\tau^{ff}=\tau^{bw}=1$, making $N^{jam}=50$
while cap\_step$=166.7$, a fundamental contradiction.
The coarse CTM spikes cell occupancies above $N^{jam}$ but drains them
in the same step, giving full capacity at a physically meaningless state.
The fine CTM with CFL-satisfying sub-steps recovers correct behaviour.
\textbf{Remedy:} always enforce $L \ge v^{ff}\,\Delta t$ at the modelling stage.

% ··············································
\subsubsection{Four-link arterial corridor}
% ··············································

Four copies of the urban arterial link (Group XXIV, $T = 120$):
\begin{center}
\begin{tabular}{lccc}
\toprule
Model & Exit vehicles & Dead-band & Prediction \\
\midrule
LTM & 1992 & 12 steps ($= 4(\tau^{ff}+1)$) & $(T-12)\times 30 = 3240$ \\
CTM & 3576 & 3 steps  ($= N_L - 1$)        & $(T-3)\times 30 = 3510$  \\
\bottomrule
\end{tabular}
\end{center}
The LTM dead-band is the sum of per-link discrete offsets; the CTM
dead-band is one step per inter-link crossing (zero-delay per link).
LTM exits 1992 vehicles $\approx 61\%$ of its ceiling (same
discretisation-induced oscillatory throughput phenomenon as Group XXII;
each link contributes a $\tau^{bw}=4$ blocking cycle).

% ··············································
\subsubsection{Travel-time accuracy: LTM offset vs CTM zero-delay}
% ··············································

Group XXV tests a 12-vehicle pulse on both a short link
($L=1\,\text{km}$, $\Delta t=1\,\text{min}$) and a freeway link
($L=10\,\text{km}$, $\Delta t=5\,\text{min}$):
\begin{center}
\begin{tabular}{lcccc}
\toprule
Link & True ff (min) & LTM 1st exit & CTM step-0 exit & LTM overestimate \\
\midrule
Short  & 1.2 & step 4 = 4 min   & 10.0 veh & $+233\%$ \\
Freeway & 6.0 & step 4 = 20 min & 10.0 veh & $+233\%$ \\
\bottomrule
\end{tabular}
\end{center}
Both models have $\tau^{ff}=2$; the LTM discrete-offset convention adds
two more steps ($\tau^{ff}+2 = 4$), placing the true delay between
the CTM zero-lower-bound and the LTM upper-bound.

% ··············································
\subsubsection{CTM vs LTM: model comparison}
\label{sec:traffic-ctm-ltm-comparison}
% ··············································

\begin{table}[htbp]
\centering
\caption{CTM vs LTM: structural and computational comparison.}
\label{tab:ctm-ltm-comparison}
\begin{tabularx}{\textwidth}{lXX}
\toprule
Feature & CTM (Daganzo 1994) & LTM (Yperman 2007) \\
\midrule
State variable &
  Cell occupancy $n_{a,m,k}$ &
  Cumulative curves $N^{in}_{a,k}$, $N^{out}_{a,k}$ \\
State size &
  $O(M_a\,|\mathcal{A}|\,T)$ &
  $O(|\mathcal{A}|\,T)$ \\
CFL &
  $\delta_a \ge v^{ff}_a\,\Delta t$ (explicit) &
  $\tau^{ff} = \lceil L/(v^{ff}\,\Delta t) \rceil$ (implicit) \\
Travel delay &
  Zero-delay artefact in single-cell implementation &
  $\tau^{ff}\,\Delta t$ in theory; current implementation adds 1--2 indexing steps \\
FIFO consistency &
  Approximate (cell-averaged) &
  Exact (cumulative-count inversion) \\
Backward wave &
  Cell-to-cell propagation &
  Exact $\tau^{bw}$ look-back on $N^{out}$ \\
Oscillatory throughput &
  Not observed in tested CFL-satisfying multi-cell cases &
  Possible under coarse time steps when discrete storage and delay are mismatched \\
Route tracking &
  Per-cell commodity fractions &
  Per-class cumulative curves \\
LP formulation size &
  $O((\sum_a M_a + |\mathcal{R}|)\,T)$ &
  $O(|\mathcal{A}|\,T)$ \\
Individual EV coupling &
  Via commodity fraction reconstruction &
  Direct FIFO inversion of $N^{out}$ \\
Typical use &
  System-level DUE, shockwave analysis &
  Meso-micro EV–charging coupling \\
Implementation &
  \texttt{ctm\_propagation.cpp} &
  \texttt{ltm\_propagation.cpp} \\
\bottomrule
\end{tabularx}
\end{table}

% ··············································
\subsubsection{Model validation tests}
\label{sec:traffic-nr-ev-penetration}

% ··············································
\subsubsection{EV penetration sweep}
% ··············································

A single station ($\bar{B}=10$ plugs, mean service $\bar{\mu}^{-1}=0.50\,$h)
on a 1-link corridor; 40\% of EVs require charging each step:
\begin{table}[htbp]
\centering
\caption{EV penetration sweep: station loading.}
\label{tab:ltm-ev-penetration}
\begin{tabular}{ccc}
\toprule
EV penetration & Max queue (veh) & Avg plug utilisation \\
\midrule
10\% &   449 & 60.3\% \\
30\% & 1\,936 & 60.6\% \\
50\% & 3\,426 & 60.6\% \\
70\% & 4\,915 & 60.6\% \\
\bottomrule
\end{tabular}
\end{table}
Queue lengths grow linearly with penetration.  At 70\% the arrival rate
($1500 \times 0.7 \times 0.4 = 420\,\text{veh/h}$) exceeds service
capacity ($10/0.5 = 20\,\text{veh/h}$) by factor 21, confirming the
runaway queue prediction.

% ··············································
\subsubsection{Individual EV trajectories and V2G}
% ··············································

Three EV agents traverse a 3-link corridor (6.5\,km total, $T=96$
steps, $\Delta t = 0.25\,$h):
\begin{table}[htbp]
\centering
\caption{Individual EV agent profiles.  Battery = 80\,kWh,
  consumption = 0.20\,kWh/km, reserve = 8\,kWh.}
\label{tab:ltm-ev-agents}
\begin{tabular}{lccc}
\toprule
Agent & Initial SOC & Initial energy (kWh) & Final SOC \\
\midrule
EV-high & 85\% & 68.0 & 83.4\% \\
EV-mid  & 45\% & 36.0 & 43.4\% \\
EV-low  & 20\% & 16.0 & 18.4\% \\
\bottomrule
\end{tabular}
\end{table}
No EV triggered the safety condition~\eqref{eq:charging-trigger-hard}.
For the most vulnerable agent (EV-low, $e_0 = 16\,\text{kWh}$), the
remaining energy after traversing the full 6.5\,km corridor is
$e_0 - E^{drive} = 16.0 - 6.5\times0.20 = 14.7\,\text{kWh}$;
since $14.7 > e^{res} = 8\,\text{kWh}$, no charging stop is triggered.
FIFO trajectory reconstruction~\eqref{eq:fifo-exit} verified for a
20-vehicle cohort: all 20 trajectories satisfy temporal consistency,
$k^{out} \ge k^{in}$, minimum travel time $\ge \tau^{ff}\,\Delta t$,
and strict FIFO order.

% ··············································
\subsubsection{Fleet assignment: station-capacity benchmark (4-node, 5 EVs)}
\label{sec:traffic-nr-fleet-assignment}
% ··············································

\paragraph{Network and parameters.}
A 4-node, 4-link network with two routes from node~1 to node~4
($\Delta t = 1\,\text{h}$, $T = 6$, default station price $0.30\,\$/\text{kWh}$).
Five EVs each request 20\,kWh of charging energy:
\begin{itemize}
  \item \textbf{Route 1} (near): links 11$\to$12; free-flow $0.50\,\text{h}$;
        station~101 at node~2 — 40\,kW total, 2 fast plugs at 20\,kW/veh.
  \item \textbf{Route 2} (far): links 21$\to$22; free-flow $1.00\,\text{h}$;
        station~102 at node~3 — 60\,kW total, 3 fast plugs at 20\,kW/veh.
\end{itemize}
All 5 vehicles have the same origin/destination (node 1$\to$4) departing at
step~0.  Table~\ref{tab:nr-fleet-assignment} compares three assignment
algorithms.

\begin{table}[htbp]
\centering
\caption{Fleet assignment comparison: 4-node benchmark, 5 EVs.
  $h_r$: route flow;
  $E^{del}$: delivered energy;
  $E^{unsrv}$: unserved energy;
  VOT~$= 1.0\,\$/\text{h}$.
  Source: \texttt{test\_ev\_power\_traffic\_report.cpp} (11 assertions, all passing).}
\label{tab:nr-fleet-assignment}
\begin{tabular}{lcccccc}
\toprule
Algorithm & $h_1$ & $h_2$ & $E^{del}$ & $E^{unsrv}$ & Travel & LP obj \\
 & (veh) & (veh) & (kWh) & (kWh) & (veh$\cdot$h) & (\$) \\
\midrule
DSP    & 5 & 0 & 40.0  & 60.0 & 2.50 & --- \\
Greedy & 2 & 3 & 100.0 & 0.0  & 4.00 & --- \\
SO-LP  & 2 & 3 & 100.0 & 0.0  & 4.00 & 34.00 \\
\bottomrule
\end{tabular}
\end{table}

\begin{remark}[SO/greedy vs.\ DUE-MSA]
  DSP, Greedy, and SO-LP are single-iteration assignment methods, not
  DUE equilibrium solvers.  They demonstrate energy delivery correctness
  and station-capacity enforcement, not Wardrop optimality.
  The CTM-MSA DUE solver is validated separately in
  Section~\ref{sec:traffic-nr-ctm-due-4node} (4-node, 4 EVs,
  \texttt{test\_ev\_power\_traffic\_ctm\_due.cpp}, 15 test cases,
  159~assertions passing), which reports final EV Wardrop gap
  $\epsilon^E$, ICV gap $\epsilon^F$, session synthesis, energy
  delivery, schedule-delay route reversal, and price-of-anarchy.
\end{remark}

\paragraph{DSP station overload.}
The Deterministic Shortest Path assigns all 5 vehicles to the near route,
ignoring station capacity.  Station~101 can supply only 40\,kW (2 plugs
$\times$ 20\,kW/plug), so 2 of 5 vehicles receive their full 20\,kWh each
($= 40$\,kWh delivered) and 3 vehicles receive nothing
($= 60$\,kWh unserved).
A queue forms: 5 vehicles are waiting in queue at step~2 (average wait
$1.75\,\text{h}$), 3 at step~3 ($1.05\,\text{h}$), 1 at step~4
($0.35\,\text{h}$), confirming the $M/D/c$-type queue dynamics of
Section~\ref{sec:traffic-charging-station}.

\paragraph{Greedy and SO-LP agreement.}
Both the Capacity-Aware Greedy and the HiGHS-certified System-Optimal LP
produce the identical split ($h_1 = 2$, $h_2 = 3$), saturating both
stations at exactly their installed capacity (40 and 60\,kW respectively)
for one step each with zero queue.
Table~\ref{tab:nr-station-step} shows the per-step station states for the
Greedy case; the SO-LP results are numerically identical.

\begin{table}[htbp]
\centering
\caption{Per-step station state --- Capacity-Aware Greedy and SO-LP (identical).
  $P^{ch}_s$: station charging power (kW);
  $B_s$: vehicles actively charging;
  $H_s$: vehicles in post-service dwell awaiting egress.
  At step~2, $B=0$ and $H>0$ because all charging completed at step~1;
  no pre-service waiting ($Q_s = 0$ throughout).}
\label{tab:nr-station-step}
\begin{tabular}{ccccccc}
\toprule
Step $k$ & $P^{ch}_{101}$ & $B_{101}$ & $H_{101}$ &
           $P^{ch}_{102}$ & $B_{102}$ & $H_{102}$ \\
 & (kW) & (veh) & (veh) & (kW) & (veh) & (veh) \\
\midrule
0 &  0 & 0 & 0 &  0 & 0 & 0 \\
1 & 40 & 2 & 0 & 60 & 3 & 0 \\
2 &  0 & 0 & 2 &  0 & 0 & 3 \\
3 &  0 & 0 & 0 &  0 & 0 & 0 \\
4 &  0 & 0 & 0 &  0 & 0 & 0 \\
5 &  0 & 0 & 0 &  0 & 0 & 0 \\
\bottomrule
\end{tabular}
\end{table}

\paragraph{System-Optimal LP objective.}
The LP minimises total system cost; with station price $0.30\,\$/\text{kWh}$
and VOT~$= 1.0\,\$/\text{h}$:
\begin{equation}
  J^{SO} =
    \underbrace{0.30 \times 100}_{\text{energy cost} = \$30.00}
    + \underbrace{1.0 \times (2\times 0.5 + 3\times 1.0)}_{\text{travel delay} = \$4.00}
    = \$34.00.
  \label{eq:nr-so-obj}
\end{equation}
The LP is proven optimal by HiGHS (NativeDualSimplex backend),
MIP gap $= 0$ (pure LP, no integer variables), primal infeasibility
$< 10^{-14}$.

% ··············································
\subsubsection{CTM-DUE MSA validation: 4-node benchmark}
\label{sec:traffic-nr-ctm-due-4node}
% ··············································

\paragraph{Setup.}
The same 4-node, 4-link symmetric network of
Section~\ref{sec:traffic-nr-fleet-assignment} is used, but solved with
the full CTM-MSA DUE algorithm (eq:ctm-due-algorithm).
4~EVs depart at step~0, each requesting 10\,kWh at station~101 (route~1)
or station~102 (route~2); $\Delta t = 1\,\text{h}$, $T = 5$,
$\epsilon^{tol} = 10^{-4}$, max 50~iterations.
Source: \texttt{test\_ev\_power\_traffic\_ctm\_due.cpp}
(15~test cases, 159~assertions, all passing; B1--B3 outputs below).

\paragraph{B1: Session synthesis and dispatch.}
\begin{center}
\begin{tabular}{lc}
\toprule
Metric & Value \\
\midrule
Converged            & yes \\
MSA iterations       & 1 \\
$\epsilon^E$ (final EV Wardrop gap, eq:ev-gap) & $0.00 \times 10^{0}$ \\
$\epsilon^F$ (final ICV Wardrop gap, eq:icv-gap) & $0.00 \times 10^{0}$ \\
Sessions synthesised & 1 \\
$E^{req}$ (total requested energy) & 40.00\,kWh \\
$E^{del}$ (total delivered energy) & 40.00\,kWh \\
$E^{unsrv}$ (unserved energy) & 0.00\,kWh \\
Service ratio        & 100.0\% \\
Route~1 flow $h_1$   & 4.000\,veh \\
Route~2 flow $h_2$   & 0.000\,veh \\
Peak station~101 load (step~1) & 42.1\,kW \\
Station~101 queue    & 0\,veh (all steps) \\
\bottomrule
\end{tabular}
\end{center}
On an uncongested symmetric network with no schedule delay all routes
have equal generalised cost at the initial equal-split, so the
relative gap $\epsilon^E = 0$ at iteration~1.  The auxiliary AoN
assignment concentrates flow on route~1 (tie-broken by index); the MSA
step produces $h_1 = 4$, $h_2 = 0$.  All 4 vehicles are served without
queueing ($W_{s,k}=0$\,hr for all $s,k$), and the single synthesised
session delivers all 40\,kWh.

\paragraph{B2: SOC feasibility filter.}
Route~1 requires $2\times0.5\,\text{km}\times2\,\text{kWh/km}=2.0\,\text{kWh}$
driving energy; route~2 requires $2\times1.5\times2=6.0\,\text{kWh}$.
With $e_0=2.0\,\text{kWh}$, route~2 is infeasible
(eq:soc-feasibility-filter) and is removed before AoN assignment:
$h_1 = 4.0000\,\text{veh}$, $h_2 = 0.000000\,\text{veh}$.
Total flow conserved: $h_1 + h_2 = 4.0$.

\paragraph{B3: Price-of-anarchy benchmark.}
$\Delta t = 0.1\,\text{h}$, $T = 6$; both DUE and SO benchmarks run.
\begin{center}
\begin{tabular}{lc}
\toprule
Metric & Value \\
\midrule
Converged              & yes \\
MSA iterations         & 1 \\
$\epsilon^E$ (final EV gap) & $0.00 \times 10^{0}$ \\
DUE TSTT               & 0.4000\,veh$\cdot$hr \\
SO TSTT                & 0.0000\,veh$\cdot$hr \\
PoA ratio              & 1.00 \\
SO benchmark solved    & yes \\
\bottomrule
\end{tabular}
\end{center}
At free-flow speed on an uncongested network the DUE TSTT equals
$4\,\text{veh}\times0.10\,\text{hr/route} = 0.40\,\text{veh}\cdot\text{hr}$.
The SO benchmark produces TSTT~$=0$ because the system-optimal LP
minimises total link flow weighted by travel time; with zero congestion
(all BPR terms at base value) the LP finds a trivial lower bound of~0.
The PoA ratio $= \text{DUE TSTT}/\text{SO TSTT}$ is reported as~1.0
when the denominator is below $10^{-9}$ (both values are near
free-flow; no congestion to exploit).

\paragraph{Schedule-delay reversal (test case~11).}
To validate eq:psi-icv and eq:icv-schedule-delay, the asymmetric network
(route~1: ff~$\approx 0.10$\,hr; route~2: ff~$\approx 0.30$\,hr) is run
with an ICV demand (6 vehicles) whose desired arrival time is
$t^* = 0.25$\,hr and early penalty $\gamma_e = 20$\,\$/hr, $\gamma_l = 0$.
Without schedule delay, route~1 wins (shorter travel time).
With the penalty:
\[
  \Psi^F_1 = 0.10 + \frac{\gamma_e \max(0,\,0.25-0.10)}{VOT}
           = 0.10 + 20\times0.15 = 3.10\,\$,
  \qquad
  \Psi^F_2 = 0.30 + 0 = 0.30\,\$.
\]
AoN assigns all 6 ICV vehicles to route~2 (verified: $h_2 > h_1$,
total $= 6.0$\,veh, $\epsilon^F \ge 0$ and finite).

\paragraph{Gap fields (test case~12).}
On a symmetric EV-only problem converged to $\epsilon^{tol}=10^{-3}$:
$\epsilon^F = 0$ (no ICV demands), $\epsilon^E \ge 0$, finite, and
$\epsilon^E < \epsilon^{tol}$ upon convergence;
\texttt{CTMDUEResult::relative\_gap}~$= \max(\epsilon^E, \epsilon^F)$
(all verified by 4 assertions).

% ··············································
\subsubsection{Joint social welfare optimizer: FormD small case (2-bus, 4 EVs)}
\label{sec:traffic-nr-formd}
% ··············································

\begin{remark}[FormD solver modes and certificates]
  The default Formulation~D solver (\texttt{joint\_optimizer.cpp}) is the
  solver-certified LP/MILP reported in this numerical table.  It uses either
  free-flow travel times, or one exogenous CTM pre-pass when
  \texttt{use\_ctm\_travel\_times=true}; in both cases the LP/MILP has a
  valid primal feasibility residual, dual bound, and MIP gap for the assembled
  linear/integer model.  A new
  \texttt{CertifiedDynamicMPECMILP} /
  \texttt{CertifiedDynamicUserBenefitMPECMILP} mode assembles a finite
  linear dynamic MPEC MILP with endogenous affine link costs, road-fault
  choice filtering, Wardrop big-\(M\) complementarity, charging/V2G, and
  optional DC-OPF in one solver call.  Its global certificate is the MILP
  certificate for that finite linear model, not a proof for nonlinear CTM/LTM.
  The newer
  \texttt{CertifiedFullJointLtmPwlMILP} /
  \texttt{CertifiedFullJointLtmUserBenefitPwlMILP} mode keeps the same
  charging/V2G, DC-OPF, and Wardrop big-\(M\) rows but adds LTM cumulative
  count variables \(N^{in}_{a,k},N^{out}_{a,k}\), LTM sending/receiving and
  availability-capacity rows, and a piecewise-linear BPR travel-cost envelope.
  Its certificate is a global MILP certificate for the assembled LTM
  PWL-MILP approximation, with \texttt{pwl\_max\_abs\_error\_bound} reporting
  the travel-cost approximation boundary.  It is explicitly not an exact
  global certificate for the original nonlinear CTM/LTM MPEC.
  A separate monolithic nonlinear mode,
  \texttt{JointOptimizerMode::FullJointSocialWelfareNLP} /
  \texttt{FullJointUserBenefitNLP}, is implemented in
  \texttt{joint\_optimizer\_full\_nlp.cpp}.  That mode assembles endogenous
  BPR congestion costs, Wardrop complementarity rows, aggregate charging/V2G
  and SOC equations, and optional DC-OPF equations in one reduced finite-route
  NLP/MPEC.  Its certificate is local NLP feasibility/stationarity only; it is
  not a global nonlinear MIP-gap proof for the full CTM/LTM MPEC.
\end{remark}

\paragraph{Network and parameters.}
A 2-bus DC power system (bus~1: generator G1 at $10\,\$/\text{MWh}$,
$\bar{p}^g_1 = 2\,\text{MW}$; bus~2: generator G2 at $50\,\$/\text{MWh}$,
$\bar{p}^g_2 = 0.5\,\text{MW}$; branch~1--2: $x = 0.1\,\text{pu}$,
rate $= 0.5\,\text{MW}$).
Traffic sub-network identical to Section~\ref{sec:traffic-nr-fleet-assignment}
(4 nodes, 2 routes).  4 EVs, WTP~$= \$50/\text{veh}$, energy request
$= 10\,\text{kWh/veh}$, VOT~$= 1.0\,\$/\text{h}$, station price
$= 0.03\,\$/\text{kWh}$; $\Delta t = 1\,\text{h}$, $T = 6$.
Station~101 on bus~1 (40\,kW, 2 plugs); station~102 on bus~2 (60\,kW,
3 plugs).

Table~\ref{tab:nr-formd} reports three LP/MILP solver modes from
\texttt{test\_ev\_power\_traffic\_joint\_opt\_d.cpp}.  The same test file also
contains a certified dynamic MPEC MILP case, local-certificate checks for the
full-joint nonlinear mode, and honesty tests that reject unsupported global
nonlinear certificates.

\begin{table}[htbp]
\centering
\caption{FormD joint optimizer: 2-bus, 4-EV results.
  SWmax = SocialWelfareMax; UBM = UserBenefitMax;
  Vars/Cons: LP variable/constraint counts;
  $W$: social welfare.}
\label{tab:nr-formd}
\begin{tabular}{lccccccc}
\toprule
Mode & $h_1$ & $h_2$ & $B_{EV}$ & $C_{del}$ & $C_{gen}$ & $W$ &
  Vars / Cons \\
 & (veh) & (veh) & (\$) & (\$) & (\$) & (\$) & \\
\midrule
SWmax, no OPF  & 4.0 & 0.0 & 200.00 & 2.00 & 0.00 & 198.00 & 11 / 27 \\
SWmax + DC-OPF & 4.0 & 0.0 & 200.00 & 2.00 & 0.40 & 197.60 & 53 / 45 \\
UBM, no OPF    & 4.0 & 0.0 & 200.00 & 2.00 & 0.00 & ---    & 11 / 27 \\
\bottomrule
\end{tabular}
\end{table}

\paragraph{Certified finite dynamic MPEC MILP numerical check.}
The D14 regression case activates
\texttt{CertifiedDynamicUserBenefitMPECMILP} with
\texttt{require\_exact\_mathematical\_model=true}.  Link~11 is unavailable at
departure step~0, so the near route is removed from the feasible
route--departure set.  The remaining far route includes a V2G-capable stop at
station~102, and bus~2 carries a small base load so reverse EV power can be
absorbed by the DC-OPF balance.  HiGHS returns an optimal MILP solution with
zero reported MIP gap, primal residual below \(10^{-6}\), and integrality
residual below \(10^{-6}\).  The result fields are
\[
\begin{gathered}
  \texttt{certified\_dynamic\_mpec\_model=true},\quad
  \texttt{travel\_times\_are\_endogenous=true},\quad
  \texttt{wardrop\_complementarity\_enforced=true},\\
  \texttt{dcopf\_coupling\_enforced=true},\quad
  \texttt{charging\_v2g\_enforced=true},\quad
  \texttt{global\_optimum\_certificate=true}.
\end{gathered}
\]
The numerical route flow is \(h_1=0\), \(h_2=4\), with all four vehicles
served.  The Wardrop complementarity residual and road-capacity violation are
both below test tolerance.  This verifies the assembled finite linear
MILP/MPEC equations; it does not re-label the nonlinear FullJoint NLP as a
global optimizer.

\paragraph{Certified full-joint LTM PWL-MILP numerical check.}
The D15 regression case activates
\texttt{CertifiedFullJointLtmUserBenefitPwlMILP} with
\texttt{require\_exact\_mathematical\_model=true}.  The case uses the same
road-failure and V2G setting as D14, but the solver model also contains
LTM cumulative-count variables, link monotonicity, sending/receiving,
entry/exit-capacity, and availability rows.  The nonlinear BPR link-cost
curve is represented by a piecewise-linear envelope with six segments; the
maximum segment error is reported through
\texttt{pwl\_max\_abs\_error\_bound}.  HiGHS returns a globally certified MILP
solution satisfying the configured MIP gap, primal residual, integrality
residual, Wardrop complementarity residual, and
\texttt{ltm\_conservation\_max\_violation} tolerances.  The asserted result
fields include
\[
\begin{gathered}
  \texttt{certified\_full\_joint\_ltm\_pwl\_milp\_model=true},\quad
  \texttt{full\_ltm\_dynamics\_enforced=true},\quad
  \texttt{pwl\_approximation\_used=true},\\
  \texttt{travel\_times\_are\_endogenous=true},\quad
  \texttt{dcopf\_coupling\_enforced=true},\quad
  \texttt{charging\_v2g\_enforced=true},\quad
  \texttt{global\_optimum\_certificate=true}.
\end{gathered}
\]
The route-flow outcome is again \(h_1=0\), \(h_2=4\), with positive V2G
energy and populated DC-OPF dispatch.  This is the strongest current
certified path for the joint model, but its honesty boundary is essential:
the certificate is for the LTM PWL-MILP approximation, not for the exact
nonlinear CTM/LTM--DUE/MPEC.

\paragraph{Full-joint nonlinear mode.}
For a finite route/departure set $\mathcal{A}_d$, the implemented monolithic
mode introduces choice flows $h_i$ and demand minimum-cost variables $\mu_d$.
Endogenous route costs are computed from BPR link costs,
\begin{equation}
  C_i(h)=\mathrm{VOT}_d\sum_{a\in r_i}
  t^0_a\left(1+\alpha_a\left(\frac{x_{a,k_i}}{\bar{x}_{a,k_i}}\right)^{\beta_a}\right)
  + C^{chg}_i + C^{sch}_i,
  \qquad
  x_{a,k}=\sum_{i:a\in r_i,\ k_i=k} h_i .
\end{equation}
Wardrop equilibrium is assembled as complementarity equations and
non-negativity inequalities:
\begin{equation}
  C_i(h)-\mu_d \ge 0,\qquad
  h_i\bigl(C_i(h)-\mu_d\bigr)=0,
  \qquad i\in\mathcal{A}_d .
\end{equation}
The implementation supplies analytic sparse derivatives to the local NLP
solver.  For link-flow sensitivity,
\begin{equation}
  \frac{\partial C_i}{\partial h_j}
  =
  \left(\mathrm{VOT}_d+\frac{\partial C^{sch}_i}{\partial t^{arr}_i}\right)
  \sum_{a\in r_i\cap r_j:\,k_i=k_j}
  t^0_a\alpha_a\beta_a
  \frac{1}{\bar{x}_{a,k_i}}
  \left(\frac{x_{a,k_i}}{\bar{x}_{a,k_i}}\right)^{\beta_a-1}.
\end{equation}
Thus the Wardrop-product row has Jacobian entries
\begin{equation}
  \frac{\partial}{\partial h_j}
  \left[h_i(C_i-\mu_d)\right]
  =
  \mathbf{1}_{i=j}(C_i-\mu_d)
  + h_i\frac{\partial C_i}{\partial h_j},
  \qquad
  \frac{\partial}{\partial \mu_d}
  \left[h_i(C_i-\mu_d)\right] = -h_i .
\end{equation}
The same NLP contains demand conservation, road-capacity inequalities,
station charge/discharge bounds, SOC recursions, terminal energy requirements,
and, when \texttt{include\_dcopf=true}, DC branch-flow and nodal power-balance
equations with EV net load
$p^{EV}_{b,k}=10^{-3}\sum_{s:\mathrm{bus}(s)=b}(P^{ch}_{s,k}-P^{dis}_{s,k})$.
These linear blocks are inserted directly into the sparse Jacobian.  When
\texttt{full\_joint\_verify\_sparse\_derivatives=true}, the code compares the
analytic gradient, equality Jacobian, and inequality Jacobian against
finite-difference values at the assembled initial point and reports
\texttt{sparse\_derivatives\_verified},
\texttt{derivative\_check\_max\_abs\_error}, and
\texttt{derivative\_check\_failures}.

The returned \texttt{JointOptimizerResult} marks
\texttt{monolithic\_full\_joint\_model=true},
\texttt{endogenous\_congestion\_enforced=true}, and
\texttt{wardrop\_complementarity\_enforced=true}.  It deliberately keeps
\texttt{global\_optimum\_certificate=false} and \texttt{proven\_optimal=false}
unless a future global MPEC/MINLP solver supplies a valid dual bound.  If
\texttt{require\_global\_nonlinear\_certificate=true}, the implementation
returns an unsupported status instead of silently treating the sparse NativeNLP
local solve as a global proof.

\paragraph{Welfare accounting.}
All four vehicles prefer the near route: WTP\,$=\$50/\text{veh}$ exceeds
both the travel delay ($0.5\,\text{h} \times \$1/\text{h} = \$0.50/\text{veh}$)
and the energy cost ($10\,\text{kWh} \times 0.03\,\$/\text{kWh} = \$0.30/\text{veh}$).
Station~101 delivers $4 \times 10 = 40\,\text{kWh}$ in one step at
$P^{ch}_{101} = 4 \times 10\,\text{kW} = 40\,\text{kW}$.

\begin{align}
  B_{EV} &= 4 \times \$50 = \$200. \\
  C_{del} &= \text{VOT} \times (4 \times 0.5\,\text{h}) = \$2.00. \\
  C_{gen} &= 0 \quad \text{(no OPF coupling in SWmax no-OPF mode)}. \\
  W &= B_{EV} - C_{del} - C_{gen} = 200 - 2 - 0 = \$198.00.
\end{align}

With DC-OPF coupling, generator G1 dispatches
$0.04\,\text{MW}$ at step~1 to balance the EV load:
\begin{equation}
  C_{gen} = 10\;\$/\text{MWh} \times 0.04\;\text{MW} \times 1\;\text{h} = \$0.40,
  \quad  W = 200 - 2 - 0.40 = \$197.60.
  \label{eq:nr-welfare-opf}
\end{equation}
The DC-OPF formulation adds 53 LP variables (generators, voltage angles,
branch flows for all $T$ steps) and 45 constraints versus the 11 variables
and 27 constraints of the traffic-only LP.

For UserBenefitMax the LP minimises $C_{del} - B_{EV} + C_{elec}$
where $C_{elec} = 4 \times 10\,\text{kWh} \times 0.03\,\$/\text{kWh} = \$1.20$,
giving objective $= 2.00 - 200.00 + 1.20 = -\$196.80$.

\paragraph{Solver certificate.}
All three modes: proven optimal, solver status \texttt{Optimal},
NativeDualSimplex backend, primal infeasibility $\le 7 \times 10^{-11}$,
integrality violation $= 0$ (pure LP).

% ··············································
\subsubsection{V2G pure-LP validation (Proposition 1)}
\label{sec:traffic-nr-v2g}
% ··············································

\paragraph{Setup.}
One EV ($n_{vehicles} = 1$) on route~1 only;
$e_0 = 50\,\text{kWh}$, $e_{min} = 20\,\text{kWh}$, $e_{max} = 60\,\text{kWh}$;
dwell $= 2$ steps; station price $\pi_s = 1.0\,\$/\text{kWh}$ (high, to
incentivise maximal discharge); $\bar{p}^{dis} = 10\,\text{kW}$,
$\bar{p}^{ch} = 20\,\text{kW}$; WTP $= \$50$; VOT $= 1.0\,\$/\text{h}$;
mode: UserBenefitMax with V2G enabled.
No binary variables are present (pure LP following
Proposition~\ref{prop:no-simultaneous}).

\begin{remark}[No-arbitrage condition in this test]
  This test does not rely on Proposition~\ref{prop:no-simultaneous} to
  enforce sequential behaviour.  The EV begins with sufficient SOC and
  the objective rewards only delivered V2G energy; charging is never
  economically attractive, so the LP solution is discharge-only.  The
  no-arbitrage condition~\eqref{eq:no-arbitrage-condition} is
  \emph{not} the reason no simultaneous charging occurs in this
  particular instance---it is a consequence of the LP cost structure
  under these specific parameters.  The no-arbitrage condition
  $\pi^{dis}\eta^{ch}\eta^{dis} \le \pi^{ch} + c^{deg}$ is a
  \emph{price-parameter condition}: when it holds, a LP solution with
  $p^{ch}_{v,k}p^{dis}_{v,k}=0$ is provably at least as good as any
  simultaneous schedule (Proposition~\ref{prop:no-simultaneous}).
  If instead $\pi^{dis}\eta^{ch}\eta^{dis} > \pi^{ch} + c^{deg}$
  (arbitrage is profitable), binary mode constraints
  $I^{ch}_{v,k}+I^{dis}_{v,k}\le 1$ would be required.
\end{remark}

\paragraph{Optimal V2G plan.}
Table~\ref{tab:nr-v2g} shows the LP-optimal per-step charging schedule.

\begin{table}[htbp]
\centering
\caption{V2G pure-LP: per-step session plan for 1 EV
  (\texttt{test\_ev\_power\_traffic\_joint\_opt\_d.cpp}, test~D9,
  all assertions passing).}
\label{tab:nr-v2g}
\begin{tabular}{ccccccc}
\toprule
Dwell step $k$ &
  $p^{ch}_{v,k}$ & $p^{dis}_{v,k}$ &
  $e_{v,k\text{-start}}$ & $e_{v,k\text{-end}}$ &
  $p^{ch}{\cdot}p^{dis}$ & Prop.~1 \\
 & (kW) & (kW) & (kWh) & (kWh) & & \\
\midrule
1 & 0.00 & 10.00 & 50 & 40 & 0 & \checkmark \\
2 & 0.00 & 10.00 & 40 & 30 & 0 & \checkmark \\
\bottomrule
\end{tabular}
\end{table}

The product $p^{ch}_{v,k} \cdot p^{dis}_{v,k} = 0$ at both steps: the EV
discharges at maximum 10\,kW for both dwell steps without any simultaneous
charging.  Under the no-arbitrage condition~\eqref{eq:no-arbitrage-condition}
(satisfied here: $\pi^{dis}\eta^{ch}\eta^{dis} = 1.0 \times 1 \times 1 \le
\pi^{ch} + c^{deg} = \infty$), the LP optimum is provably non-simultaneous
by Proposition~\ref{prop:no-simultaneous}, so no binary mode constraints
are required \emph{for this test instance}.  If the no-arbitrage condition
were violated, binary constraints $I^{ch}_{v,k}+I^{dis}_{v,k}\le 1$
would be needed.

\paragraph{Objective decomposition.}
Total V2G energy $= 2 \times 10\,\text{kW} \times 1\,\text{h} = 20\,\text{kWh}$;
V2G revenue $= 20\,\text{kWh} \times \$1.0/\text{kWh} = \$20$;
travel delay $C_{del} = 1.0 \times 0.5\,\text{h} = \$0.50$:
\begin{equation}
  J^{UBM}
  = C_{del} - B_{EV} - E^{dis}_{tot} \cdot \pi_s
  = 0.50 - 50.00 - 20.00 \cdot 1.0
  = -\$69.50.
  \label{eq:nr-v2g-obj}
\end{equation}
Verified by test assertion: \texttt{res.objective} $= -69.50$
(tolerance $10^{-6}$); \texttt{proven\_optimal = true};
primal infeasibility $< 10^{-7}$.

\paragraph{Simulation cross-check.}
Test Scenario~4 in \texttt{test\_ev\_power\_traffic\_report.cpp}
independently simulates one V2G session ($e_0 = 50\,\text{kWh}$,
$p^{dis} = 20\,\text{kW}$, 1 dwell step): final SOC $= 30\,\text{kWh}$,
$E^{dis}_{tot} = 20\,\text{kWh}$, EV power balance margin increases by
$0.020\,\text{MW}$ at the V2G step.
Both the simulation test and the joint-optimizer LP agree on
$E^{dis} = 20\,\text{kWh}$.

% ··············································
\subsubsection{Large-scale benchmark: 10$\times$10 grid, 10,000 EVs, IEEE-118}
\label{sec:traffic-nr-largescale}
% ··············································

\paragraph{System description.}
\begin{itemize}
  \item \textbf{Power grid}: IEEE 118-bus; buses 1--100 each host a
        fast-charging station with 300 fast plugs at 50\,kW each
        (15\,MW/station, \$0.30/kWh).
  \item \textbf{Traffic grid}: 10$\times$10 Cartesian grid, 100 nodes,
        180 directed links; free-flow $0.25\,\text{h/link}$,
        capacity $2000\,\text{veh/h}$.
  \item \textbf{Corridors}: 5 north--south corridors, each with two routes.
        Route A: 9 vertical links down the left column, 1 horizontal link
        at the bottom (charging stop at mid-left node, row~4).
        Route B: 1 horizontal link at the top, 9 vertical links down the
        right column (charging stop at mid-right node, row~4).
        $\Rightarrow$ 10 routes total.
  \item \textbf{Demands}: 5 corridors $\times$ 4 departure steps $= 20$
        fleet demands, 500 vehicles/demand $= 10\,000$ vehicles total.
        Each vehicle requests 30\,kWh at up to 50\,kW.
  \item $\Delta t = 1\,\text{h}$, $T = 12$;
        enforce road capacity and station capacity.
\end{itemize}

\paragraph{Algorithm comparison.}
Table~\ref{tab:nr-largescale} reports results from
\texttt{test\_ev\_power\_traffic\_largescale.cpp} (14 assertions, all passing).

\begin{remark}[Scalability vs.\ DUE validation]
  The algorithms compared here (CapacityAwareGreedy and SystemOptimalLP)
  are assignment heuristics and a system-optimal LP, respectively; they
  are not DUE-MSA runs.  This benchmark demonstrates scalability
  (10,000 EVs, 100 stations, IEEE-118 power grid) and
  energy-delivery correctness, not equilibrium convergence.
  A DUE-MSA run on this network would additionally report:
  final ICV gap $\epsilon^F$, final EV gap $\epsilon^E$, number of MSA
  iterations, departure-time split, and used-path cost dispersion.
  These metrics are omitted here because the test uses a single-iteration
  greedy/SO-LP approach for performance benchmarking.
\end{remark}

\begin{table}[htbp]
\centering
\caption{Large-scale benchmark: 10$\times$10 grid, 10,000 EVs, IEEE-118.
  All 10,000 vehicles served, $E^{del} = 300{,}000\,\text{kWh}$,
  $E^{unsrv} = 0$.}
\label{tab:nr-largescale}
\begin{tabular}{lccccc}
\toprule
Algorithm & Solve time & Travel time & Peak station & LP obj & Proved \\
 & (ms) & (veh$\cdot$h) & (kW) & & optimal \\
\midrule
CapacityAwareGreedy & 0.5  & 25{,}001 & 9{,}000  & ---       & no  \\
SystemOptimalLP     & 3.1  & 25{,}015 & 15{,}000 & 115{,}000 & yes \\
\bottomrule
\end{tabular}
\end{table}

\paragraph{Route assignment.}
Under CapacityAwareGreedy each of the 20 demands splits 300/200
(Route~A/Route~B), distributing load between mid-left stations
(41, 43, 45, 47, 49) and mid-right stations (42, 44, 46, 48, 50).
Under SystemOptimalLP all 500 vehicles per demand concentrate on
Route~A (LP objective $= 115{,}000$), loading left-column stations
at their 15\,MW installed capacity but achieving certified global
optimality.

\paragraph{Station utilisation (CapacityAwareGreedy).}
Table~\ref{tab:nr-station-util} reports peak load and total delivered
energy for the 10 route stations across all 4 departure steps.

\begin{table}[htbp]
\centering
\caption{Per-station utilisation under CapacityAwareGreedy.
  Route-A stations (odd IDs) serve 300 vehicles/demand;
  Route-B stations (even IDs) serve 200 vehicles/demand;
  each station active for 4 departure-step groups.}
\label{tab:nr-station-util}
\begin{tabular}{lccc}
\toprule
Station (bus) & Route & Peak load (kW) & Delivered energy (kWh) \\
\midrule
41 & A, corridor 0 & 9{,}000 & 36{,}000 \\
42 & B, corridor 0 & 6{,}000 & 24{,}000 \\
43 & A, corridor 1 & 9{,}000 & 36{,}000 \\
44 & B, corridor 1 & 6{,}000 & 24{,}000 \\
45 & A, corridor 2 & 9{,}000 & 36{,}000 \\
46 & B, corridor 2 & 6{,}000 & 24{,}000 \\
47 & A, corridor 3 & 9{,}000 & 36{,}000 \\
48 & B, corridor 3 & 6{,}000 & 24{,}000 \\
49 & A, corridor 4 & 9{,}000 & 36{,}000 \\
50 & B, corridor 4 & 6{,}000 & 24{,}000 \\
\midrule
Total & & --- & 300{,}000 \\
\bottomrule
\end{tabular}
\end{table}

Energy balance verified: $10{,}000 \times 30\,\text{kWh} = 300{,}000\,\text{kWh}$,
matching \texttt{total\_delivered\_energy\_kwh} to within 1\,kWh.
Route-A peak: each vehicle requests $30\,\text{kWh}$ in a 1-hour step,
requiring an average charging power of $30\,\text{kWh}/1\,\text{h} = 30\,\text{kW}$
(within the 50\,kW plug hardware limit);
$300\,\text{veh} \times 30\,\text{kW/veh} = 9{,}000\,\text{kW}$.
Route-B analogously: $200 \times 30 = 6{,}000\,\text{kW}$.
Under SystemOptimalLP all 500 vehicles in each demand window are concentrated
on Route~A ($500 \times 30 = 15{,}000\,\text{kW}$, exactly saturating the
15\,MW station capacity).

\paragraph{Scalability.}
Both solvers complete well within their declared time limits
(CapacityAwareGreedy: 0.5\,ms vs.\ 5\,s limit;
SystemOptimalLP: 3.1\,ms vs.\ 30\,s limit), demonstrating that
the LP and greedy heuristic both scale efficiently to 10,000-vehicle
fleets on a 100-node grid.

% ------------------------------------------------------------
\subsection{Extension Interfaces to Power and Resilience Models}
\label{sec:traffic-extension}
% ------------------------------------------------------------

\begin{center}
\begin{tabularx}{\textwidth}{L{0.32\textwidth}X}
\toprule
\textbf{Output quantity} & \textbf{Consuming module} \\
\midrule
Station charging power $p^{ch}_{s,k}$
  & Smart-charging / DC-OPF
    (Section~\ref{sec:ev_power_traffic_dynamic_model}) \\
Station dwell population $H_{s,k}$ and per-EV SOC
  & V2G aggregator $\to$ OPF \\
EV departure flow $o_{s,k}$
  & LTM downstream receiving update \\
Aggregate load $P^{EV}_{b,k}$
  & Power-network load profile \\
V2G potential $P^{v2g,pot}_{s,k}$
  & Reserve and resilience models \\
\bottomrule
\end{tabularx}
\end{center}

% ============================================================
\subsection{Optimization Framework}
\label{sec:traffic-opt}
% ============================================================

The simulation sub-modules use MSA for DUE and rule-based heuristics for
charging admission.  This section gives proper mathematical programmes that
replace or wrap those heuristic layers, providing optimality guarantees and
enabling joint power-traffic problems.

\begin{remark}[Implementation status of optimisation sub-modules]
  Table~\ref{tab:cpp-extension} lists the source files.
  The default joint social welfare optimiser is a solver-certified LP/MILP
  over route, charging/V2G, per-step road-capacity, and optional DC-OPF rows.
  Its travel times are exogenous: either free-flow values or a one-shot
  equal-split CTM pre-pass.  The separate FullJoint NLP mode is the only
  current path with endogenous congestion rows, and it is locally certified
  only.  The remaining five formulations compile, expose direct
  public entry points, and are covered by
  \texttt{test\_ev\_power\_traffic\_formulations\_e\_f\_g\_h}.  In the
  second implementation round, the optimisation modules were upgraded from
  static prototypes to solver-assembled dynamic programmes with explicit
  certificates:
  \begin{itemize}
    \item FormD full-joint NLP: a new monolithic reduced NLP/MPEC mode
          assembles endogenous BPR congestion, Wardrop complementarity,
          charging/V2G, and optional DC-OPF in one programme.  It reports
          local feasibility/stationarity, uses analytic sparse derivatives,
          and explicitly refuses a requested global nonlinear certificate
          because the current in-tree solvers do not prove a global
          MPEC/MINLP gap for this model.
    \item SO-CTM LP: route-specific cell states, inter-cell flows, explicit
          node movement variables $y_{a,b,k}$, and aggregate merge/origin
          receiving constraints are assembled in the LP.
    \item SO-LTM LP: per-route cumulative counts are connected by explicit
          node movement variables; downstream receiving constraints aggregate
          all incoming movements and origin injection at the same link-step.
    \item DUE / Frank-Wolfe: full endogenous EV departure-window assignment
          is available through \texttt{DUEMode::FullEndogenous}; the final
          solution reports disaggregated departure-route flows, VI residual,
          normalised VI gap, complementarity residual, and demand-conservation
          residual computed at final CTM costs.
    \item Infrastructure MILP: BPR Beckmann piecewise-linear cost is
          available; an optional time-expanded LTM dynamic envelope enforces
          route-link cumulative counts, node movement, capacity, and
          availability constraints inside the MILP.
    \item LTM-MPC: per-station access-link $N_{\text{in}}/N_{\text{out}}$
          tracking with strict SOC--V2G linkage; access-link inflows are
          bounded by an exogenous OD/route forecast so the controller cannot
          create artificial arrivals.
  \end{itemize}
  The remaining distinction is not merely cosmetic: DUE fixed points are
  checked by residual metrics, LP/MILP modules are certified only for their
  assembled linear/integer subproblems, and the nonlinear FullJoint mode has no
  global MPEC/MINLP proof.
  For paper-audit runs the C++ options expose
  \texttt{require\_exact\_mathematical\_model}.  When this flag is set, every
  reduced, residual-based, forecast-driven, aggregate, or locally certified
  path returns an \texttt{unsupported} status rather than a numerical result.
  Therefore a reported \texttt{mathematical\_model\_verified=true} cannot be
  produced by the present implementation unless a future solver path supplies
  an equation-by-equation full-model certificate.
\end{remark}

\subsubsection{CTM-based system-optimal traffic assignment (SO-CTM LP)}
\label{sec:traffic-ctm-so-lp}

Under minimisation of total vehicle-time the non-smooth $\min(\cdot,\cdot)$
in~\eqref{eq:ctm-intercell-flow} is linearised by two one-sided upper bounds
(Ziliaskopoulos 2000).

\paragraph{Objective.}
\begin{equation}
  \min\; \sum_{a,m,k} n_{a,m,k}\,\Delta t.
  \label{eq:so-ctm-obj}
\end{equation}

\paragraph{Constraints.}
\begin{align}
  & n_{a,m,k+1} = n_{a,m,k} + y_{a,m-1,k} - y_{a,m,k}
    && \forall\, a,m,k
    \label{eq:so-ctm-balance}\\
  & y_{a,m,k} \le \tfrac{v^{ff}_a}{\delta_a}\,n_{a,m,k}\,\Delta t
    && \forall\, a,m,k
    \label{eq:so-ctm-send-vf}\\
  & y_{a,m,k} \le \bar{q}_a\,\Delta t
    && \forall\, a,m,k\\
  & y_{a,m,k} \le \tfrac{w_a}{\delta_a}
                  (N^{jam}_{a,m+1} - n_{a,m+1,k})\,\Delta t
    && \forall\, a,\, m < M_a-1,\, k
    \label{eq:so-ctm-recv}\\
  & \textstyle\sum_{r \in \mathcal{R}(od)} x_{r,k} = d^{od}_k
    && \forall\, od, k\\
  & n_{a,m,0} = 0
    && \forall\, a,m
    \label{eq:so-ctm-ic}
\end{align}
The initial empty condition $n_{a,m,0}=0$ is standard.  The terminal
clearance constraint $n_{a,m,T}=0$ is valid only when the planning
horizon $T$ is long enough that all demand has been served; for shorter
horizons, replace it with a terminal penalty $\lambda_T\sum_{a,m}n_{a,m,T}$
in the objective~\eqref{eq:so-ctm-obj}.

The receiving constraint~\eqref{eq:so-ctm-recv} applies to internal
cell-to-cell flows ($m < M_a-1$).  For the exit flow from the last
cell $m = M_a-1$, the downstream receiving capacity is imposed by the
node model of the successor link or sink: if link $a$ discharges to a
sink, the receiving capacity is set to the sink capacity (or $\infty$);
for an ordinary downstream link the first-cell receiving constraint of
that link provides the bound via the CTM node equations.
LP size: $O((\sum_a M_a + |\mathcal{R}|)\,T)$ variables,
$O(\sum_a M_a\,T)$ constraints.

\subsubsection{LTM-based system-optimal assignment (SO-LTM LP)}
\label{sec:traffic-ltm-so-lp}

\paragraph{Objective.}
\begin{equation}
  \min\; \sum_{a,k} \bigl(N^{in}_{a,k} - N^{out}_{a,k}\bigr)\,\Delta t
  \;=\;
  \min\; \sum_{a,k} n_{a,k}\,\Delta t,
  \label{eq:so-ltm-obj}
\end{equation}
the total vehicle-hours in the network (minimising total delay).
The terminal residual $\sum_a(N^{in}_{a,T}-N^{out}_{a,T})$ may be used
as an alternative when the planning horizon $T$ is long enough that the
network empties; it is easier to evaluate but ignores travel time
throughout the simulation period.

\paragraph{Constraints (all linear).}
\begin{align}
  & N^{out}_{a,k} - N^{out}_{a,k-1}
      \le N^{in}_{a,k-\tau^{ff}_a} - N^{out}_{a,k-1}
    && \text{(sending)}
    \label{eq:so-ltm-send}\\
  & N^{out}_{a,k} - N^{out}_{a,k-1} \le \bar{q}_a\,\Delta t
    && \text{(cap)}\\
  & N^{in}_{a,k} - N^{in}_{a,k-1}
      \le N^{out}_{a,k-\tau^{bw}_a} + N^{jam}_a - N^{in}_{a,k-1}
    && \text{(receiving)}
    \label{eq:so-ltm-recv}\\
  & N^{in}_{a,k} - N^{in}_{a,k-1} \le \bar{q}_a\,\Delta t
    && \text{(cap)}\\
  & \textstyle\sum_b y_{a,b,k} = N^{out}_{a,k} - N^{out}_{a,k-1}
    && \text{(exit split)}\\
  & \textstyle\sum_a y_{a,b,k} + d^{orig}_{b,k}
      = N^{in}_{b,k} - N^{in}_{b,k-1}
    && \text{(entry merge)}\\
  & N^{in}_{a,k} \ge N^{in}_{a,k-1},\quad
    N^{out}_{a,k} \ge N^{out}_{a,k-1}
    && \text{(monotone)}
\end{align}
LP size: $O(|\mathcal{A}|\,T)$, independent of cell count.

\subsubsection{DUE as a variational inequality}
\label{sec:traffic-due-vi}

Let the feasible departure-rate set be
\begin{equation}
  \mathcal{H}
  = \Bigl\{
    (h^F, h^E) \ge 0 :
    \sum_{k \in \mathcal{K}^{dep}_w}\sum_{r} h^F_{w,r,k} = D^F_w,\;
    \sum_{k \in \mathcal{K}^{dep}_w}\sum_{\omega} h^E_{w,\omega,k} = D^E_w,\;
    \forall\, w
  \Bigr\}.
  \label{eq:due-feasible-set}
\end{equation}
The full-DUE feasible set $\mathcal{H}$ from~\eqref{eq:due-feasible-set}
contrasts with the fixed-departure feasible set:
\begin{equation}
  \mathcal{H}^{fix}
  = \Bigl\{
    (h^F, h^E) \ge 0 :
    \sum_{r} h^F_{w,r,k} = d^F_{w,k},\;
    \sum_{\omega} h^E_{w,\omega,k} = d^E_{w,k},\;
    \forall\, w, k
  \Bigr\}.
  \label{eq:due-feasible-set-fix}
\end{equation}
The mixed-fleet DUE condition~\eqref{eq:icv-wardrop}--\eqref{eq:ev-wardrop-zero}
is equivalent to the variational inequality: find
$(h^{F,*}, h^{E,*}) \in \mathcal{H}$ such that
\begin{equation}
  \sum_{w,r,k}
    \Psi^F_{w,r,k}(h^*)
    \bigl(h^F_{w,r,k} - h^{F,*}_{w,r,k}\bigr)
  +
  \sum_{w,\omega,k}
    \Psi^E_{w,\omega,k}(h^*)
    \bigl(h^E_{w,\omega,k} - h^{E,*}_{w,\omega,k}\bigr)
  \ge 0
  \quad
  \forall\; (h^F, h^E) \in \mathcal{H}.
  \label{eq:due-vi}
\end{equation}
For static or separable link-cost approximations the VI can be reduced
to the Beckmann potential
\begin{equation}
  \min_{(h^F,h^E) \in \mathcal{H}}
  \Phi(h) = \sum_a \int_0^{f_a(h)} t_a(u)\,\mathrm{d}u
  \label{eq:beckmann}
\end{equation}
which is convex when link cost functions are monotone and separable, and
can be solved by Frank-Wolfe or interior-point methods.  In the full
dynamic CTM/LTM setting with spillback and station queues, path costs
$\Psi^F, \Psi^E$ depend on departure time and link costs are
non-separable; the DUE is therefore treated as the fixed-point
problem~\eqref{eq:fixed-point} or the VI~\eqref{eq:due-vi} rather than
a Beckmann program.  The Beckmann formulation and its uniqueness
certificate are available only under the separable-monotone static
approximation.

\subsubsection{Joint power-traffic social welfare optimisation}
\label{sec:traffic-joint-sw}

\begin{equation}
  \max_{x,\,p^g,\,p^{ch}}\;
    \sum_{d} W^P_d\,s_d
    - \sum_g C_g(p^g)
    - \mathrm{VOT} \sum_r x_r\,T^r(x),
  \label{eq:joint-sw}
\end{equation}
subject to DC power balance, LTM traffic dynamics, and charging-load
coupling.  The general coupling is expressed through station net load:
\begin{equation}
  p^{EV}_{b,k}
  = \frac{1}{1000}
    \sum_{s:\mathrm{bus}(s)=b}
    \bigl(P^{ch}_{s,k} - P^{dis}_{s,k}\bigr)
    \quad\text{[MW]},
  \label{eq:joint-ev-load}
\end{equation}
where
\begin{equation}
  P^{ch}_{s,k}
  = \sum_{\omega\ni s} x^E_{\omega,k}\,p^{ch}_{\omega,k},
  \qquad
  P^{dis}_{s,k}
  = \sum_{\omega\ni s} x^E_{\omega,k}\,p^{dis}_{\omega,k}.
\end{equation}
\begin{remark}[$P^{EV}$ sign convention]
  In this traffic module, positive $p^{EV}_{b,k}$ denotes \emph{net EV
  load} (aggregate charging exceeds discharging).  In power-flow equations
  using injection-positive convention, EV load enters as $-p^{EV}_{b,k}$,
  so that generation equals load plus net EV demand.
\end{remark}
In simplified benchmark cases without V2G and with all admitted vehicles
charging at a fixed plug power, this reduces to
\begin{equation}
  p^{EV}_{b,k}
  = \frac{1}{1000}
    \sum_{s:\mathrm{bus}(s)=b}
    \sum_{r\ni s} x_{r,k}\,\bar{p}^{ch}_s.
\end{equation}
The current benchmark tests use this simplified form; the station
net-load form~\eqref{eq:joint-ev-load} is the mathematical interface
for full V2G coupling.

\begin{remark}[Implemented solver: LP/MILP, certified finite MILP/MPEC, and reduced full-joint NLP modes]
  The C++ implementation in \texttt{joint\_optimizer.cpp}
  (\texttt{solve\_joint\_optimizer}) supports the solver-certified LP/MILP
  modes described above.  In the default static mode, free-flow route travel
  times $T^{FF}_r$ are used; when
  \texttt{JointOptimizerOptions::use\_ctm\_travel\_times=true}, one CTM
  forward pass supplies exogenous $T^{CTM}_{r,k_{dep}}$ before the LP/MILP is
  assembled.  The
  \texttt{CertifiedDynamicMPECMILP} and
  \texttt{CertifiedDynamicUserBenefitMPECMILP} modes add a finite linear
  dynamic MPEC MILP with affine endogenous congestion, availability/fault
  filtering, Wardrop big-\(M\) rows, charging/V2G, and optional DC-OPF; its
  global certificate is the MILP certificate for that finite model.  In
  addition, \texttt{joint\_optimizer\_full\_nlp.cpp} implements
  a reduced monolithic nonlinear NLP/MPEC with endogenous BPR congestion,
  charging/V2G, DC-OPF, and Wardrop complementarity in one programme.  This
  nonlinear mode is an exact assembly of the finite-route reduced equations,
  with analytic sparse derivatives and optional finite-difference derivative
  verification, but the current in-tree solve path reports only a local NLP
  feasibility/stationarity certificate.  It is not a global certificate for
  the full CTM/LTM nonlinear MPEC.
\end{remark}

\subsubsection{MILP for charging infrastructure design}
\label{sec:traffic-infra-milp}

Let $z_s \in \{0,1\}$ and $\bar{B}_s \in \mathbb{Z}_+$:
\begin{equation}
  \min_{z,\bar{B},x,n}\;
    \sum_s \bigl(c^{cap}_s\,\bar{B}_s\,z_s
               + c^{op}_s \sum_k p^{ch}_{s,k}\,\Delta t\bigr)
    + \mathrm{VOT} \sum_r x_r\,T^r,
  \label{eq:infra-milp-obj}
\end{equation}
with $B_{s,k} \le \bar{B}_s\,z_s$ and $\bar{B}_s \le \bar{B}^{max}_s\,z_s$
(big-$M$ linking constraints).  The implemented solver uses static OD
conservation, station power/plug constraints, optional BPR link-flow variables,
and native branch-and-cut through \texttt{engine::solve\_milp\_bc}.  When
\texttt{InfraDesignMILPOptions::include\_ltm\_dynamic\_constraints} is enabled,
the MILP also includes time-expanded LTM cumulative-count variables, node
movement variables, link capacity/availability constraints, and downstream
merge receiving constraints.  This gives a solver-certified dynamic design
envelope; a full CTM cell-state infrastructure MILP is intentionally left out
because its variable count grows with $\sum_a M_aT$.

\subsubsection{LTM-based model predictive control (MPC)}
\label{sec:traffic-mpc}

At each step $k$, solve a finite-horizon LP over horizon $H$ with stage cost
\begin{equation}
  \ell_{k+j} = \omega_Q \sum_s Q_{s,k+j}
             + \omega_T \sum_a (N^{in}_{a,k+j} - N^{out}_{a,k+j})
             - \omega_{V2G} \sum_s P^{v2g,pot}_{s,k+j},
  \label{eq:mpc-stage-cost}
\end{equation}
controlled via admission rates $b_{s,k}, \ldots, b_{s,k+H-1}$.
Only the first-step action $b_{s,k}$ is applied; the LP is resolved
at step $k+1$ with the observed state.

\paragraph{Recommended integration sequence.}
\begin{enumerate}
  \item \textbf{Offline design:} SO-CTM LP or infrastructure MILP to
        select station locations, capacities, and base route flows.
  \item \textbf{Day-ahead:} joint social welfare optimisation
        (\texttt{solve\_joint\_social\_welfare}) to set station prices
        $\pi_{s,k}$ aligning EV charging with grid marginal costs.
  \item \textbf{Real-time:} LTM-MPC adjusting admission rates in
        response to observed queue lengths and grid emergency signals.
\end{enumerate}

\begin{table}[htbp]
\centering
\caption{C++ extension map for optimisation-based traffic models.
  \textbf{Implemented}: source file compiles, links, and has a direct
  regression test.
  \textbf{Implemented}$^\dagger$: source exists but approximates the exact
  mathematical formulation described in this chapter.}
\label{tab:cpp-extension}
\begin{tabularx}{\textwidth}{lXXl}
\toprule
Formulation & Source file & Notes & Status \\
\midrule
Joint SW optimizer &
  \texttt{joint\_optimizer.cpp} /
  \texttt{assignment\_lp.cpp} &
  Free-flow (default) or CTM-dynamic travel times
  (\texttt{use\_ctm\_travel\_times}); per-step
  road-capacity bounds; HiGHS LP; DC-OPF optional &
  \textbf{Implemented} \\
SO-CTM LP &
  \texttt{ctm\_so\_lp.cpp} &
  \texttt{engine::LPModel}, HiGHS;
  multi-commodity per-route cell states $n^{(r)}_{a,m,k}$ with
  aggregate capacity constraints (equations \eqref{eq:so-ctm-balance}--\eqref{eq:so-ctm-recv}
  use the aggregate single-commodity $n_{a,m,k}$);
  flexible origin injection; link-availability support;
  explicit node movement variables and aggregate downstream receiving
  constraints; $M=\max(1, \texttt{n\_cells\_per\_link})$ is
  user-supplied and is not auto-enforced for CFL stability &
  \textbf{Implemented}$^\dagger$ \\
SO-LTM LP &
  \texttt{ltm\_so\_lp.cpp} &
  \texttt{engine::LPModel}, HiGHS;
  multi-commodity per-route cumulative counts
  $N^{\text{in},(r)}/N^{\text{out},(r)}$ with aggregate node
  capacity/receiving constraints
  (equations \eqref{eq:so-ltm-send}--\eqref{eq:so-ltm-recv}
  use link-aggregate variables);
  correct origin injection; link-availability support;
  explicit $y_{a,b,k}$ movement variables (per-route \texttt{Ynode})
  and merge receiving constraints;
  $\texttt{so\_tstt\_hr}=\text{LP objective}$ &
  \textbf{Implemented}$^\dagger$ \\
DUE / Frank-Wolfe &
  \texttt{ctm\_due\_vi.cpp} &
  Frank-Wolfe with CTM forward pass; EV-demand only
  (eq.~\eqref{eq:due-vi} includes mixed ICV+EV terms $\Psi^F+\Psi^E$;
  ICV traffic handled separately via \texttt{simulate\_icv\_traffic});
  golden-section line search; departure-window support;
  Wardrop/VI residual stopping metrics; reports VI gap,
  normalised VI gap, complementarity residual, and demand-conservation
  residual (fixed-point certificates, not LP dual bounds) &
  \textbf{Implemented}$^\dagger$ \\
Infrastructure MILP &
  \texttt{infra\_design\_milp.cpp} &
  \texttt{engine::MIPModel}, branch-and-cut;
  binary $z_s$ and integer $\bar{B}_s$;
  BPR Beckmann piecewise-linear congestion cost for the
  objective (first-order derivative approximation at $v=\bar{q}_a$;
  not the exact $\mathrm{VOT}\!\sum_r x_r T^r$ from
  eq.~\eqref{eq:infra-milp-obj});
  optional time-expanded LTM dynamic envelope (Nin/Nout/Ynode with
  zero objective coefficient — constraints only, no VHT
  minimisation when the LTM layer is active) &
  \textbf{Implemented}$^\dagger$ \\
LTM-MPC &
  \texttt{ltm\_mpc.cpp} &
  Receding-horizon admission LP per station;
  LTM $N_{\text{in}}/N_{\text{out}}$ tracking on per-station
  \emph{access links only} (eq.~\eqref{eq:mpc-stage-cost}
  sums $\omega_T$ over all network links $a$; the code covers
  access links only, which is a proxy for full-network VHT);
  strict SOC--V2G dwell linkage;
  OD-forecast arrival bound on access-link inflow
  (static free-flow proportional allocation) &
  \textbf{Implemented}$^\dagger$ \\
\bottomrule
\end{tabularx}
\end{table}

\begin{remark}[Residual gaps between code and mathematical model]
\label{rem:implementation-gaps}
All five formulations have direct regression tests
(\texttt{test\_ev\_power\_traffic\_formulations\_e\_f\_g\_h.cpp}) that
verify LP/MILP feasibility, primal-certificate residuals, node-flow
conservation, and demand conservation at the solver-returned solution.
The following gaps between the \emph{exact} mathematical model and the
current implementation remain:

\begin{itemize}
  \item \textbf{SO-CTM LP and SO-LTM LP --- aggregate vs.\ per-route
    state representation.}
    \texttt{CTMOptions::use\_aggregate\_link\_variables} (and the
    corresponding \texttt{LTMSOLPOptions} flag) switches the LP to reduced
    \emph{aggregate} link/cell state variables $n_{a,m,k}$ and
    $N^{in/out}_{a,k}$ shared across routes.  Per-route entry/exit or
    node-movement variables are still retained for OD accounting, so this is
    an aggregate relaxation/approximation of the full route-commodity model,
    not a proof that route-level FIFO is exactly enforced.  When the flag is
    \texttt{false} (default), the per-route multi-commodity decomposition is
    used instead.  Result fields
    \texttt{per\_route\_decomposition\_used},
    \texttt{closed\_link\_handling\_strict}, and the LTM aggregate
    strictness flags expose which scope was actually solved.

  \item \textbf{SO-CTM LP --- CFL stability auto-enforced (resolved).}
    The Courant--Friedrichs--Lewy ceiling
    $M^{\mathrm{CFL}}_{a} = \lfloor \ell_a / (v^{ff}_a\,\Delta t) \rfloor$
    is now computed per link and the user-supplied
    $\mathtt{CTMOptions::n\_cells\_per\_link}$ is clamped to
    $\min(M_{\mathrm{user}}, M^{\mathrm{CFL}}_a)$ at setup time.
    When the user supplies an explicit $\Delta t_{\mathrm{CTM}} > 0$ the
    clamp applies to that value; otherwise
    $\Delta t_{\mathrm{CTM}}$ is auto-computed from
    $\mathtt{choose\_ctm\_dt}$ with a $0.9\,v^{ff}$ safety margin
    that already satisfies CFL.

  \item \textbf{DUE-VI --- mixed-fleet certificate scope.}
    \texttt{CTMDUEVIOptions::include\_icv\_in\_due} enables ICV
    participation in the Frank--Wolfe loop.  When set, both EV and ICV
    route flows are updated via auxiliary-assignment steps
    (\texttt{auxiliary\_assignment\_icv}) and convergence is checked using
    the worst EV/ICV Wardrop gap available at the final CTM costs.  However,
    \texttt{compute\_due\_vi\_certificate} still computes the VI residual
    over EV demand variables only.  Therefore, when ICVs participate,
    \texttt{CTMDUEResult::vi\_certificate\_covers\_ev\_only=true}; the ICV
    side is reported through \texttt{icv\_relative\_gap}, not a full
    mixed-fleet VI dual certificate.

  \item \textbf{Infrastructure MILP --- BPR objective approximation.}
    The objective~\eqref{eq:infra-milp-obj} specifies
    $\mathrm{VOT}\sum_r x_r T^r(x)$ with endogenous travel times.
    The code replaces this with a piecewise-linear BPR Beckmann
    approximation; the excess-flow slope $\alpha\beta=0.6$ matches
    the first-order derivative at $v_a=\bar{q}_a$.  When
    \texttt{include\_ltm\_dynamic\_constraints} is enabled, the
    Nin/Nout variables receive objective coefficients
    $+\mathrm{VOT}\cdot\Delta t$ and $-\mathrm{VOT}\cdot\Delta t$
    respectively, so their difference $N^{in}_{a,k}-N^{out}_{a,k}$
    proxies vehicle-hours in transit; the previously double-counted
    static route-cost $h$-variable and BPR link-flow terms are zeroed
    out when the LTM block is active.  Ynode movement-flow variables
    carry zero cost (they are kinematic feasibility variables).

  \item \textbf{Certified dynamic MPEC MILP --- global certificate scope.}
    The \texttt{CertifiedDynamic*MPECMILP} modes are globally certified only
    for their finite linear model: affine link-delay costs, admitted
    route--departure choices, big-\(M\) Wardrop rows, charging/V2G, and
    optional DC-OPF.  They deliberately do not claim exact nonlinear CTM cell
    propagation, LTM FIFO reconstruction, nonlinear BPR powers, or a global
    certificate for \texttt{FullJoint*NLP}.  The Wardrop cost uses
    route-level planned charging/V2G energy terms; optimized session powers
    are coupled through linear SOC and station rows, but their realized
    power-per-served-vehicle ratios are not placed back into the Wardrop cost,
    because that would create bilinear equations outside the certified MILP
    class.  Therefore,
    \texttt{mathematical\_model\_verified=true} in these modes means
    ``verified for the assembled finite MILP/MPEC'', not ``the whole
    paper-level nonlinear CTM/LTM MPEC has been globally solved''.

  \item \textbf{LTM-MPC --- tracked-link VHT proxy.}
    \texttt{LTMMPCOptions::track\_all\_route\_links} extends the MPC LP
    with $N^{in}_{a,k}/N^{out}_{a,k}$ variables for \emph{every} link
    on each station's approach route.  The $\omega_T$ stage cost then
    sums over all tracked links, which is closer to
    stage cost~\eqref{eq:mpc-stage-cost} than the single-access-link proxy.
    The admission constraint
    remains anchored to the final access link only.  When the flag is
    \texttt{false} (default), only the single access-link proxy is used.
    The OD-forecast arrival bound continues to use a static free-flow
    proportional allocation
    (\texttt{forecast\_station\_arrivals\_from\_demands}), not a
    dynamic traffic prediction.
\end{itemize}
\end{remark}
